\section{Derived Functors}

\subsection{$\delta$-functors}

Recall that morphisms of exact sequences are collections of morphisms which commute with everything.
That is to say, if we have two exact sequences $\set{A_n,f_n}_n$ and $\set{B_n,g_n}_n$ (i.e.\ $f_n\colon A_n\to A_{n+1}$ and $\im f_n=\ker f_{n+1}$), a morphism $A\to B$
is a collection $\set{s_n\colon A_n\to B_n}_n$ such that $sf=gs$ (i.e.\ view the exact sequences as chain complexes).

\bdefn

    A {\emphcolor homological} (resp.\ {\emphcolor cohomological}) {\emphcolor $\delta$-functor} from $\A$ to $\B$ (Abelian categories) is a collection of
    additive functors $T_n\colon\A\to\B$ (resp.\ $T^n\colon\A\to\B$) for $n\geq0$, along with morphisms $\delta_n\colon T_n(C)\to T_{n-1}(A)$ (resp.\ 
    $\delta^n\colon T^n(C)\to T^{n+1}(A)$) for every short exact sequence $0\to A\to B\to C\to0$ in $\A$.

    A $\delta$-functor must satisfy the following (set $T_n=T^n=0$ for $n<0$):
    \benum
        \item for a short exact sequence $0\to A\to B\to C\to0$, we get a long exact sequence
        $$ \cdots\to T_{n+1}(C) \xtos\delta T_n(A) \to T_n(B) \to T_n(C) \xtos\delta T_{n-1}(A) \to \cdots $$
        (ascending indices for cohomological $\delta$-functors.)
        In particular $T_0$ is right exact (since $T_{-1}=0$) and $T^0$ is left exact ($T_n,T^n$ in general are not: we have an exact sequence $T_n(A)\to T_n(B)\to T_n(C)$ but we cannot append $0$).

        \item For a morphism of short exact sequences, $\delta$ commutes with the induced morphisms.
        That is, if we have a morphism between short exact sequences $0\to A'\to B'\to C'\to0$ and $0\to A\to B\to C\to0$, we get commutative squares

        \bigskip
        \centerline{\drawdiagram{
            $T_n(C')$&$T_{n-1}(A')$\cr
            $T_n(C)$&$T_{n-1}(A)$
        }{
            \diagarrow{from={1,1}, to={1,2}, text=$\delta$, y distance=.25cm}
            \diagarrow{from={2,1}, to={2,2}, text=$\delta$, y distance=.25cm}
            \diagarrow{from={1,1}, to={2,1}}
            \diagarrow{from={1,2}, to={2,2}}
        }}

        for every $n\geq0$.
    \eenum

\edefn

\bprop

    Homology gives a homological $\delta$-functor $H_\bul\colon\Ch_{\geq0}(\A)\to\A$; cohomology gives a cohomological $\delta$-functor $H^\bul\colon\Ch^{\geq0}(\A)\to\A$.

\eprop

\bproof

    The $\delta$ morphisms are the connecting maps $\partial$.
    The axioms are satisfied because the connecting maps are functorial.
    \qed

\eproof

\bprop

    Let $\S(\A)$ be the category of short exact sequences in $\A$.
    If $T$ is a $\delta$-functor then $\delta_n$ is a natural transformation from $T_n$ to $T_{n-1}$ (viewed as functors mapping $0\to A\to B\to C\to0$ to $T_nC$ and $T_{n-1}A$ respectively).

\eprop

\bproof

    This is immediate from the second axiom of $\delta$-functors.
    \qed

\eproof

\bexam

    Let $p\in\bZ$ then define a homological $\delta$-functor $\Ab\to\Ab$ by $T_0(A)=A/pA$ and $T_1(A)={}_pA=\set{a\in A}[pa=0]$ ($T_n=0$ for all other $n$).
    To get the $\delta$ morphisms, the snake lemma on the diagram

    \centerline{\drawdiagram{
        $0$&$A$&$B$&$C$&$0$\cr
        $0$&$A$&$B$&$C$&$0$
    }{
        \diagarrow{from={1,1}, to={1,2}}
        \diagarrow{from={1,2}, to={1,3}}
        \diagarrow{from={1,3}, to={1,4}}
        \diagarrow{from={1,4}, to={1,5}}
        \diagarrow{from={2,1}, to={2,2}}
        \diagarrow{from={2,2}, to={2,3}}
        \diagarrow{from={2,3}, to={2,4}}
        \diagarrow{from={2,4}, to={2,5}}
        \diagarrow{from={1,2}, to={2,2}, text=$p$, x distance=.25cm}
        \diagarrow{from={1,3}, to={2,3}, text=$p$, x distance=.25cm}
        \diagarrow{from={1,4}, to={2,4}, text=$p$, x distance=.25cm}
    }}

    Since $\ker p={}_pA$ and $\coker p=A/pA$ we get an exact sequence
    $$ 0 \to {}_pA \to {}_pB \to {}_pC \xtos\delta A/pA \to B/pB \to C/pC \to 0 $$
    Which is precisely the first axiom.
    Functorality of the connecting morphism gives us the second axiom.

    We can also view this as a cohomological $\delta$-functor by swapping $T_0$ and $T_1$.

\eexam

\bdefn

    A morphism $S\to T$ of $\delta$-functors is a system of natural transformations $S_n\to T_n$ (or $S^n\to T^n$) which commute with $\delta$.

    A homological $\delta$-functor $T$ is {\emphcolor universal} if given any other homological $\delta$-functor $S$ and a natural transformation $f_0\colon S_0\to T_0$,
    we can uniquely extend $f_0$ to a morphism $\set{f_n}_n\colon S\to T$.
    Dually a cohomological $\delta$-functor $T$ is {\emphcolor universal} if given any other cohomological $\delta$-functor $S$ and natural $f_0\colon T^0\to S^0$ it can be uniquely
    extended to $T\to S$.

\edefn

\bexam

    We will show later that homology and cohomology are universal.

\eexam

\bprop

    If $F\colon\A\to\B$ is exact, then $T_0=F$ and $T_n=0$ for $n\neq0$ is a universal $\delta$-functor (both homological and cohomological).

\eprop

\bproof

    A quick check of the axioms shows that $F$ is indeed a $\delta$-functor (with $\delta=0$).
    Let $S$ be another $\delta$-functor (suppose homological) and if $f_0\colon S_0\to F$ is natural then define $f_n\colon S_n\to T_n=0$ to be zero.
    This is the only way to extend $f_0$, and it is a morphism.
    \qed

\eproof

\subsection{Projective resolutions}

\bdefn

    An object $P$ in an Abelian category $\A$ is {\emphcolor projective} if it satisfies the following universal lifting property: given an epimorphism $g\colon B\to C$
    and a map $\gamma\colon P\to C$ there exists a (not necessarily unique) map $\beta\colon P\to B$ such that $\gamma=g\circ\beta$.
    In a diagram:

    \centerline{\drawdiagram{
        &$P$\cr
        $B$&$C$&$0$
    }{
        \diagarrow{from={2,1}, to={2,2}}
        \diagarrow{from={2,2}, to={2,3}}
        \diagarrow{from={1,2}, to={2,2}, text=$\gamma$, x distance=.25cm}
        \diagarrow{from={1,2}, to={2,1}, text=$\exists\beta$, x distance=-.25cm, dashed}
    }}

\edefn

We are mostly concerned with {\emphcolor projective modules} (projective objects in $\Mod_R$).
Clearly free modules are projective: take a basis $(v_i)_i$ and map $v_i$ to a preimage of $\gamma(v_i)$.

\bprop

    Direct summands of projective modules are projective.

\eprop

\bproof

    If $P=M\oplus M'$ is projective then any map $\gamma\colon M\to C$ can be extended to a map $\gamma\colon P\to C$ and the rest follows.
    \qed

\eproof

\bprop

    An $R$-module is projective iff it is a direct summand of a free $R$-module (that is, $M$ is projective iff there exists $M'$ such that $M\oplus M'$ is free).

\eprop

\bproof

    Direct summands of a free module are by the previous proposition projective.

    Conversely $F(A)$ be the free $R$-module on the set $A$, for a module $M$ we can consider $F(M)$ on the set underlying $M$ and we have a surjection $\pi\colon F(M)\to M$.
    If $M$ is projective then the universal lifting property gives us a map $i\colon M\to F(M)$ such that $\pi i=\id_M$:

    \centerline{\drawdiagram{
        &$M$\cr
        $F(M)$&$M$&$0$
    }{
        \diagarrow{from={2,1}, to={2,2}, text=$\pi$, y distance=.25cm}
        \diagarrow{from={2,2}, to={2,3}}
        \diagarrow{from={1,2}, to={2,2}, text=$\id_M$, x distance=.25cm}
        \diagarrow{from={1,2}, to={2,1}, text=$\exists i$, x distance=-.25cm, dashed}
    }}

    So $i\pi$ is a projection onto $i(M)\cong M$, meaning $M$ is a direct summand of $F(M)$.
    \qed

\eproof

Over many nice rings (e.g.\ $\bZ$, division rings), projective modules are all free.
This is not always the case:

\bexam

    If $R=R_1\times R_2$ is the product of two rings then $R_1\times0,0\times R_2$ are projective as their sum is $R$ (a free $R$-module).
    $P=R_1\times0$ is not free since $(0,1)P=0$ (meaning for any choice $B\subseteq P$ there is an $r\in R$ such that $rB=0$ so $B$ cannot be linearly independent).
    So for example $\bZ/6\bZ\cong\bZ/2\bZ\times\bZ/3\bZ$ is projective but not free.

\eexam

\bexam

    Let $k$ be a field and consider the matrix ring $R=M_n(k)$.
    This has a module $k^n$ (multiplication is given by vector-matrix multiplication), and as an $R$-module we have $R\cong k^n\oplus\cdots\oplus k^n$, so $k^n$ is projective.
    A free $R$-module needs to have dimension $dn^2$ over $k$ for some $d>0$, but $\dim_kk^n=n$ so it cannot be free.

\eexam

The category of finite Abelian groups has no projective objects (other than $0$): non-trivial subgroups of free Abelian groups are torsion-free and thus infinite.

\bdefn

    $\A$ has {\emphcolor enough projectives} if for every $A\in\A$ there is an epimorphism $P\to A$ with $P$ projective.

\edefn

$\Mod_R$ has enough projectives, take the epimorphism $F(M)\to M$.

\blemm

    $M$ is projective iff $\hom_A(M,\bul)$ is exact.

\elemm

\bproof

    Suppose $\hom(M,\bul)$ is exact, and we have a surjection $g\colon B\to C$ and a map $\gamma\colon M\to C$.
    Since $g$ is surjective and $\hom(M,\bul)$ is exact, $g_*$ is surjective as well.
    Thus we can lift $\gamma$ to $\beta$ such that $\gamma=g_*\beta=g\circ\beta$; $M$ has the universal lifting property, it is projective.

    Recall that in general $\hom(M,\bul)$ is left exact.
    So to show that it is exact, we must only show that for $0\to A\to B\xto gC\to0$, $g_*$ is surjective as well (i.e.\ we can append $\to0$).
    This is precisely what the universal lifting property says.
    \qed

\eproof

\bdefn

    A chain complex where every component is projective is called a {\emphcolor chain complex of projectives}.
    It is not necessarily projective in $\Ch$.

\edefn

\bprop

    A chain complex $P$ is projective in $\Ch$ iff it is a split exact complex of projectives.

\eprop

\bproof

    If $P$ is projective then so is $P[-1]$: a surjection $B\to C$ and a map $P[-1]\to C$ can be shifted up to $B[1]\to C[1]$ and $P\to C[1]$, then the resulting
    lift can be shifted back down $P[-1]\to B$ ($[p]$ is functorial).
    We have a short exact sequence
    $$ 0 \to P \to \cone(P) \to P[-1] \to 0 $$
    Now a surjection onto a projective object gives rise to a projection: we can lift $\id_{P[-1]}$ to $\beta\colon P[-1]\to\cone(P)$ such that $\delta\beta=\id_{P[-1]}$ where $\delta$ is the surjection.
    Then $\beta\delta$ is a projection onto $P[-1]$ whose kernel is $P$, thus $\cone(P)\cong P\oplus P[-1]$.
    Since $\cone(P)$ is exact, so too must $P$ be.

    Since $\delta(p,q)=-p$ we have that $\beta$ is of the form $\beta(p)=(-p,s(p))$ for $s\colon P[-1]\to P$ (not necessarily a chain map).
    Now we claim that $s$ is a splitting map.
    Since $\beta$ is a chain map we have
    $$ d\beta(p) = d(-p,s(p)) = (dp,ds(p)+p),\qquad \beta(-dp) = (dp,-s(dp)) $$
    So we have $ds+\id=-s d$ so $\id=-sd-ds$ meaning $P$ is split exact.

    Now we must show that $P$'s components are all projectives.
    So we need to show that each $B'_n$ is projective: let $g\colon N\to M$ surjective and $\gamma\colon B'_n\to M$.
    Considering $N$ and $M$ as chain complexes centered at $n$ with all other terms $0$, $g$ gives rise to a surjection of a chain complex, and $\gamma$ a map from $P$ to $M$.
    Then we can lift it to $\beta\colon P\to N$, and we get the desired result.

    Recall that if $P$ is split exact then $P_n\cong B_n\oplus B'_n$ for $B'_n\cong B_{n-1}$.
    The differential is zero on $B_n$ and an isomorphism $B'_n\cong B_{n-1}$, so we can view $P_n\cong B_n\oplus B_{n-1}$ with $d(b,b')=(b',0)$.
    Now a map $\gamma\colon P\to C$ is of the form $(\gamma_n,\gamma'_n)\colon P_n\to C_n$ and this is a chain map iff $d\gamma'_n=\gamma_{n-1}$.
    Indeed: $(\gamma_{n-1},\gamma'_{n-1})d(b,b')=\gamma_{n-1}(b')$ and $d(\gamma_n(b)+\gamma'_n(b'))=d\gamma_nb+d\gamma'_nb'$.
    This is true for all $b,b'$ and so take $b=0$ and we get $\gamma_{n-1}=d\gamma'_n$.
    And conversely this implies $d\gamma_n=0$ and so we have equality.

    So if $P$ is split exact of projectives, let $(\gamma_n,\gamma'_n)\colon P\to C$ and $g\colon A\to C$ surjective.
    Each $B_n$ is projective (as it is a direct summand of a projective), so we have maps $\beta_n\colon B_{n-1}\to A_n$ which lift $\gamma'_n$ along $g_n$.
    Notice then that the map $d\beta_{n+1}\colon B_n\to A_n$ lifts $\gamma_n$ along $g_n$:
    $$ g_nd\beta_{n+1} = dg_{n+1}\beta_{n+1} = d\gamma'_{n+1} = \gamma_n $$
    Then the map $(d\beta_{n+1},\beta_n)\colon P\to A$ is a chain map and it lifts $\gamma+\gamma'$ along $g$ as desired.
    \qed

\eproof

\bcoro

    If $\A$ has enough injectives, so does $\Ch(\A)$.

\ecoro

\bproof

    Let $C$ be a chain complex, then for each $C_n$ let $g_n\colon B_n\to C_n$ be a surjection from a projective module.
    Then define $P_n=B_n\oplus B_{n-1}$ with differential mapping $d(b,b')=(b',0)$.
    This is a split exact complex of projectives, thus a projective complex.
    Now $(dg_{n+1},g_n)\colon P\to C$ is a chain map as discussed in the previous proof, and is surjective as desired.
    \qed

\eproof

\bdefn

    A short exact sequence $0\to A\to B\to C\to0$ is called {\emphcolor split} if there is an isomorphism $f\colon B\to A\oplus C$ that commutes with the short
    exact sequence (with inclusion $\iota\colon A\to A\oplus C$ and projection $\pi\colon A\oplus C\to C$).
    Explicitly, the following diagram commutes:

    \centerline{\drawdiagram{
        $0$&$A$&$B$&$C$&$0$\cr
        &&$A\oplus C$
    }{
        \diagarrow{from={1,1}, to={1,2}}
        \diagarrow{from={1,2}, to={1,3}}
        \diagarrow{from={1,3}, to={1,4}}
        \diagarrow{from={1,4}, to={1,5}}
        \diagarrow{from={1,2}, to={2,3}, text=$\iota$, x distance=-.25cm}
        \diagarrow{from={1,3}, to={2,3}, text=$f$, x distance=-.25cm}
        \diagarrow{from={2,3}, to={1,4}, text=$\pi$, x distance=-.25cm}
    }}

\edefn

\blemm[title=Splitting lemma]

    For a split exact sequence $0\to A\xto fB\xto gC\to0$, the following are equivalent:
    \benum
        \item the sequence splits;
        \item the sequence is {\emphcolor left split}: there exists a morphism $t\colon B\to A$ such that $tf=\id_A$;
        \item the sequence is {\emphcolor right split}: there exists a morphism $s\colon C\to B$ such that $gs=\id_C$.
    \eenum

\elemm

\bproof

    The sequence being split implies being left and right split: define $t\colon A\oplus C\cong B\to A$ to project onto $A$ and $s\colon C\to B\cong A\oplus C$ to
    embed $C$ into $B$.
    If the sequence is left split then $B=\ker t\oplus\im f$ since their intersection is zero and $b=(b-ft(b))+ft(b)$.
    Now $\im f\cong A$ and $\ker t\cong C$ since we can map $b\in\ker t$ to $g(b)\in C$: this map is injective since its kernel is $\ker t\cap\ker g=\ker t\cap\im f=0$.
    It is surjective since $g(b-ft(b))=g(b)$.
    Thus $B\cong A\oplus C$ and this commutes with the sequence.
    \qed

\eproof

Note that this lemma is not true in the category of groups: being (left) split is equivalent to $B$ being the direct product of $A$ and $C$, but being right split
is equivalent to $B$ being the semidirect product of $A$ and $C$.
Thus (left) split implies right split, but not the other way around.

\bprop

    An object $P\in\A$ is projective iff every short exact sequence $0\to A\to B\to P\to0$ is split.

\eprop

\bproof

    If $P$ is projective then we can lift $\id_P$ along the surjection $B\to P$ to give us a map which composes with it to give the identity.
    Thus the sequence is right split, and thus split.

    Conversely suppose every such short exact sequence is split.
    Now $\A$ is finitely complete and thus for a surjection $A\to B$ and a map $P\to B$ we can consider the diagram's pullback $C$:

    \centerline{\drawdiagram{
        $C$&$P$&$0$\cr
        $A$&$B$&$0$\cr
    }{
        \diagarrow{from={1,1}, to={1,2}, text=$f$, y distance=-.25cm}
        \diagarrow{from={1,2}, to={1,3}}
        \diagarrow{from={2,1}, to={2,2}}
        \diagarrow{from={2,2}, to={2,3}}
        \diagarrow{from={1,1}, to={2,1}}
        \diagarrow{from={1,2}, to={2,2}}
    }}

    It can be shown that the first row is indeed exact ($f$ is epi), since in Abelian categories epimorphisms are regular (cokernels of their kernels), and we can follow
    the kernel of the bottom map.
    Taking $K$ to be the kernel of $f$, we have an exact sequence $0\to K\to C\xto fP\to0$, which must split.
    Then there is a map $g\colon P\to C$ such that $fg=\id_P$.
    Now consider $P\xto gC\to A$: by commutativity we have that it gives the desired lifting of $P\to B$.
    \qed

\eproof

\bdefn

    Let $M\in\A$, a {\emphcolor left resolution} of $M$ is an exact sequence
    $$ \cdots\xtos d P_1\xtos dP_0\xtos\epsilon M\to0 $$
    The map $\epsilon$ is called the {\emphcolor augmentation map}.
    This is a {\emphcolor projective resolution} if each $P_i$ are projective.

\edefn

\blemm

    If $\A$ has enough projectives, every object $M$ has a projective resolution.
    In particular every $R$-module has a projective resolution.

\elemm

\bproof

    First take a projective $P_0$ and a surjection $\epsilon_0\colon P_0\to M$, and set $M_0=\ker\epsilon_0$.
    Inductively given $M_{n-1}$ choose a projective $P_n$ and surjection $\epsilon_n\colon P_n\to M_{n-1}$.
    Define $M_n=\ker\epsilon_n$ and let $d_n$ be the composite $P_n\to M_{n-1}\to P_{n-1}$.
    Since $M_{n-1}\to P_{n-1}$ is monic $\ker d_n=\ker\epsilon_n=M_n$.
    And we have $d_n(P_n)=M_{n-1}=\ker d_{n-1}$, thus the resulting sequence is exact.
    \qed

\eproof

\bprop

    A left resolution $\cdots\to P_1\to P_0\to M\to0$ is the same as a quasi-isomorphism $P\to M$ where $P$ is the chain complex $\cdots\to P_1\to P_0\to0$ and $M$
    is considered as a chain complex focused at $0$.

\eprop

\bproof

    A chain map $P\xtos\epsilon M$ is given by a map $P_0\xto\epsilon M$ such that $\epsilon d=0$ (for $d\colon P_1\to P_0$).
    This is a quasi-isomorphism if $\epsilon\colon H_0(P)=P_0/dP_1\to M$ is an isomorphism, meaning $\ker\epsilon=\im d$, i.e.\ the sequence is exact.
    \qed

\eproof

\bthrm[title=Comparison theorem]

    Let $P\xtos\epsilon M$ be a projective resolution of $M$ and $f'\colon M\to N$ a map in $\A$.
    Then for every resolution $Q\xtos\eta N$, there is a chain map $f\colon P\to Q$ lifting $f'$ in the sense that $\eta f_0=f'\circ\epsilon$:

    \centerline{\drawdiagram{
        $\cdots$&$P_1$&$P_0$&$M$&$0$\cr
        $\cdots$&$Q_1$&$Q_0$&$N$&$0$\cr
    }{
        \diagarrow{from={1,1}, to={1,2}}
        \diagarrow{from={1,2}, to={1,3}}
        \diagarrow{from={1,3}, to={1,4}}
        \diagarrow{from={1,4}, to={1,5}}
        \diagarrow{from={2,1}, to={2,2}}
        \diagarrow{from={2,2}, to={2,3}}
        \diagarrow{from={2,3}, to={2,4}}
        \diagarrow{from={2,4}, to={2,5}}
        \diagarrow{from={1,2}, to={2,2}}
        \diagarrow{from={1,3}, to={2,3}}
        \diagarrow{from={1,4}, to={2,4}, text=$f'$, x distance=-.25cm}
    }}

    $f$ is unique up to chain homotopy equivalence.

\ethrm

Alternatively we could say that there exists a unique (up to chain homotopy equivalence) chain map between the chain complexes $P\to M$ to $Q\to N$.

The proof will show that $P\to M$ being a projective resolution is stronger a condition than necessary: it is sufficient to be given a chain complex $P\to M\to0$
where $P_n$ are all projective.

\bproof

    We construct $f_n$ inductively; let $f_{-1}=f'$ and $f_{-2}=0$.
    Suppose $f_i$ has been constructed for $i\leq n$, so that $f_{i-1}d=df_i$.
    To construct $f_{n+1}$ we consider the $n$-cycles of $P,Q$: if $n=-1$ then set $Z_{-1}(P)=M$ and $Z_{-1}(Q)=N$.
    Since $f_{n-1}d=df_n$ we know that $f_n$ induces a map $f'_n\colon Z_n(P)\to Z_n(Q)$, so we have two diagrams with exact rows:

    \hbox{\hfill\drawdiagram{
        $\cdots$&$P_{n+1}$&$Z_n(P)$&$0$\cr
        $\cdots$&$Q_{n+1}$&$Z_n(Q)$&$0$\cr
    }{
        \diagarrow{from={1,1}, to={1,2}}
        \diagarrow{from={1,2}, to={1,3}}
        \diagarrow{from={1,3}, to={1,4}}
        \diagarrow{from={2,1}, to={2,2}}
        \diagarrow{from={2,2}, to={2,3}}
        \diagarrow{from={2,3}, to={2,4}}
        \diagarrow{from={1,2}, to={2,2}, dashed}
        \diagarrow{from={1,3}, to={2,3}, text=$f'_n$, x distance=.25cm}
    }%
    \hfill\drawdiagram{
        $0$&$Z_n(P)$&$P_n$&$P_{n-1}$\cr
        $0$&$Z_n(Q)$&$Q_n$&$Q_{n-1}$\cr
    }{
        \diagarrow{from={1,1}, to={1,2}}
        \diagarrow{from={1,2}, to={1,3}}
        \diagarrow{from={1,3}, to={1,4}}
        \diagarrow{from={2,1}, to={2,2}}
        \diagarrow{from={2,2}, to={2,3}}
        \diagarrow{from={2,3}, to={2,4}}
        \diagarrow{from={1,2}, to={2,2}, text=$f'_n$, x distance=.25cm}
        \diagarrow{from={1,3}, to={2,3}, text=$f_n$, x distance=.25cm}
        \diagarrow{from={1,4}, to={2,4}, text=$f_{n-1}$, x distance=.25cm}
    }\hfill}

    Using the first diagram and $P_{n+1}$ being projective allows us to lift $f'_nd$ to $f_{n+1}\colon P_{n+1}\to Q_{n+1}$ such that $df_{n+1}=f'_nd=f_nd$.
    And so $\set{f_n}$ is a chain map extending $f'$ as desired.

    To show uniquess, suppose $g\colon P\to Q$ is a chain map extending $f'$ as well.
    Let $h=f-g$ and we will construct a contraction $\set{s_n\colon P_n\to Q_{n+1}}$ inductively.
    For negative $n$ we must have $s_n=0$.
    For $n=0$ we have $\eta h_0=(f'-f')\epsilon=0$, and so $h_0$ maps $P_0$ to $\ker\eta=Z_0(Q)=dQ_1$.
    And so we can lift $h_0$ to $s_0\colon P_0\to Q_1$ such that $h_0=ds_0=ds_0+s_{-1}d$.

    Inductively suppose for $i<n$ we have maps $s_i$ such that $ds_{n-1}=h_{n-1}-s_{n-2}d$ and then consider $h_n-s_{n-1}d\colon P_n\to Q_n$.
    Note
    $$ d(h_n-s_{n-1}d) = dh_n - ds_{n-1}d = dh_n - (h_{n-1} - s_{n-2}d)d = dh_n - h_{n-1}d = 0 $$
    Thus $h_n-s_{n-1}d$ maps into $Z_n(Q)=dQ_{n+1}$, and so we can lift it to a map $s_{n+1}\colon P_n\to Q_{n+1}$ such that $ds_{n+1}=h_n-s_{n-1}d$.
    And so we have $h=ds+sd$ as desired.
    \qed

\eproof

Note that we did not use $P\to M$'s exactness: only $Q\to M$'s.

\blemm[title=Horeshoe lemma]

    Suppose we have the following diagram:

    \centerline{\drawdiagram{
        &&&$0$\cr
        $\cdots$&$P'_1$&$P'_0$&$A'$&$0$\cr
        &&&$A$\cr
        $\cdots$&$P''_1$&$P''_0$&$A''$&$0$\cr
        &&&$0$\cr
    }{
        \diagarrow{from={2,1}, to={2,2}}
        \diagarrow{from={2,2}, to={2,3}}
        \diagarrow{from={2,3}, to={2,4}, text=$\epsilon'$, y distance=.25cm}
        \diagarrow{from={2,4}, to={2,5}}
        %
        \diagarrow{from={4,1}, to={4,2}}
        \diagarrow{from={4,2}, to={4,3}}
        \diagarrow{from={4,3}, to={4,4}, text=$\epsilon''$, y distance=.25cm}
        \diagarrow{from={4,4}, to={4,5}}
        %
        \diagarrow{from={1,4}, to={2,4}}
        \diagarrow{from={2,4}, to={3,4}, text=$i_A$, x distance=.25cm}
        \diagarrow{from={3,4}, to={4,4}, text=$\pi_A$, x distance=.25cm}
        \diagarrow{from={4,4}, to={5,4}}
    }}

    such that the column is exact and the rows are projective resolutions.
    Set $P_n=P'_n\oplus P''_n$, then $P_n$ form a projective resolution $P$ of $A$, and the colimn lifts to an exact sequence of complexes $0\to P'\xto iP\xto\pi P''\to0$,
    where $i_n\colon P'_n\to P_n$ and $\pi_n\colon P_n\to P''_n$ are the inclusion and projection maps.

\elemm

\bproof

    Lift $\epsilon''$ to a map $P''_0\to A$; then the direct sum of this map with $i_A\epsilon'\colon P_0'\to A$ gives a map $\epsilon\colon P_0\to A$.
    Then the following diagram commutes:

    \centerline{\drawdiagram{
        &$0$&$0$&$0$\cr
        $0$&$\ker\epsilon'$&$P_0'$&$A'$&$0$\cr
        $0$&$\ker\epsilon$&$P_0$&$A$&$0$\cr
        $0$&$\ker\epsilon''$&$P_0''$&$A''$&$0$\cr
        &$0$&$0$&$0$\cr
    }{
        \diagarrow{from={2,1}, to={2,2}}
        \diagarrow{from={2,2}, to={2,3}}
        \diagarrow{from={2,3}, to={2,4}, text=$\epsilon'$, y distance=.25cm}
        \diagarrow{from={2,4}, to={2,5}}
        \diagarrow{from={3,1}, to={3,2}}
        \diagarrow{from={3,2}, to={3,3}}
        \diagarrow{from={3,3}, to={3,4}, text=$\epsilon$, y distance=.25cm}
        \diagarrow{from={3,4}, to={3,5}}
        \diagarrow{from={4,1}, to={4,2}}
        \diagarrow{from={4,2}, to={4,3}}
        \diagarrow{from={4,3}, to={4,4}, text=$\epsilon''$, y distance=.25cm}
        \diagarrow{from={4,4}, to={4,5}}
        \diagarrow{from={1,2}, to={2,2}}
        \diagarrow{from={2,2}, to={3,2}}
        \diagarrow{from={3,2}, to={4,2}}
        \diagarrow{from={4,2}, to={5,2}}
        \diagarrow{from={1,3}, to={2,3}}
        \diagarrow{from={2,3}, to={3,3}}
        \diagarrow{from={3,3}, to={4,3}}
        \diagarrow{from={4,3}, to={5,3}}
        \diagarrow{from={1,4}, to={2,4}}
        \diagarrow{from={2,4}, to={3,4}}
        \diagarrow{from={3,4}, to={4,4}}
        \diagarrow{from={4,4}, to={5,4}}
    }}

    The two right columns are exact, and all the top and bottom rows are as well.
    The snake lemma gives us that the left column is exact, and $\coker\epsilon=0$ (since $0=\coker\epsilon'\to\coker\epsilon\to\coker\epsilon''=0$ is exact).
    Thus $\epsilon$ is epic, as desired.

    Now we have another horseshoe:

    \centerline{\drawdiagram{
        &&&$0$\cr
        $\cdots$&$P'_2$&$P'_1$&$\ker\epsilon'$&$0$\cr
        &&&$\ker\epsilon$\cr
        $\cdots$&$P''_2$&$P''_1$&$\ker\epsilon''$&$0$\cr
        &&&$0$\cr
    }{
        \diagarrow{from={2,1}, to={2,2}}
        \diagarrow{from={2,2}, to={2,3}}
        \diagarrow{from={2,3}, to={2,4}, text=$d'$, y distance=.25cm}
        \diagarrow{from={2,4}, to={2,5}}
        %
        \diagarrow{from={4,1}, to={4,2}}
        \diagarrow{from={4,2}, to={4,3}}
        \diagarrow{from={4,3}, to={4,4}, text=$d''$, y distance=.25cm}
        \diagarrow{from={4,4}, to={4,5}}
        %
        \diagarrow{from={1,4}, to={2,4}}
        \diagarrow{from={2,4}, to={3,4}}
        \diagarrow{from={3,4}, to={4,4}}
        \diagarrow{from={4,4}, to={5,4}}
    }}

    and we can induct.
    \qed

\eproof

\subsection{Injective resolutions}

We can duallize the notion of projective objects:

\bdefn

    An object $I$ in an Abelian category $\A$ is {\emphcolor injective} if it satisfies the following universal lifting property:
    given a monomorphism $f\colon A\to B$ and a map $\alpha\colon A\to I$, we can lift $\alpha$ to a map $\beta\colon B\to I$ such that $\alpha=\beta\circ f$.

    \centerline{\drawdiagram{
        $0$&$A$&$B$\cr
        &$I$\cr
    }{
        \diagarrow{from={1,1}, to={1,2}}
        \diagarrow{from={1,2}, to={1,3}, text=$f$, y distance=-.25cm}
        \diagarrow{from={1,2}, to={2,2}, text=$\alpha$, x distance=-.25cm}
        \diagarrow{from={1,3}, to={2,2}, text=$\beta$, x distance=.25cm, dashed}
    }}

    $\A$ has enough injectives if for every object $A\in\A$ there is an injective $I\in\A$ and a monomorphism $A\to I$.

\edefn

Note that if $\set{I_\alpha}$ are injectives, so is their product (if it exists).

\blemm

    An injective object in $\A$ is precisely a projective object in $\A^\op$.
    $\A$ has enough injectives iff $\A^\op$ has enough projectives.

\elemm

\bcoro

    $I$ is injective iff $\hom_\A(\bul,I)$ is exact.

\ecoro

\bproof

    $I\in\A$ is injective iff $I\in\A^\op$ is projective iff $\hom_{A^\op}(I,\bul)=\hom_A(\bul,I)$ is exact.
    \qed

\eproof

\bprop[title=Baer's criterion]

    A right $R$-module $E$ is injective iff for every right ideal $J$ of $R$, every map $J\to E$ can extend to a map $R\to E$.

\eprop

\bproof

    If $E$ is injective then we lift $J\to E$ along the injection $J\to R$.
    Conversely suppose we have an $R$-module $B$ and a submodule $A$ along with a map $\alpha\colon A\to E$.
    Let $\c E$ be the poset of all extensions of $\alpha$: maps $\alpha'\colon A'\to E$ for $A\subseteq A'\subseteq B$, partially ordered by $\alpha'\leq\alpha''$ iff $\alpha''$ extends $\alpha'$.
    This clearly satisfies the criteria for Zorn's lemma, as such it has a maximal element, $\alpha'\colon A'\to E$.
    We claim that $A'=B$.

    Suppose not; there is some $b\in B$ not in $A'$.
    Define then $J=\set{r\in R}[br\in A']$: a right ideal of $R$.
    Then the composition $J\xtos bA'\xtos{\alpha'}E$ extends to a map $f\colon R\to E$.
    Let $A''=A'+bR$, and define $\alpha''\colon A''\to E$ by $\alpha''(a+br)=\alpha'(a)+f(r)$.
    This is well-defined since for $br\in A'$ we know $\alpha'(br)=f(r)$ by definition.
    As such $\alpha''$ is a proper extension of $\alpha'$, a contradiction.
    \qed

\eproof

Recall that a semisimple $R$-module is one for which every submodule has a complement: if $A\leq M$ is a submodule then there exists $A'\leq M$ such that $M=A\oplus A'$.
A semisimple ring is one which is semisimple as a module over itself (every ideal has a complement).
Clearly then modules over semisimple rings are injective: let $J'$ be $J$'s complement and extend $J\to E$ to $R=J\oplus J'\to E$.

\bcoro

    Let $R$ be semisimple, then every $R$-module is injective.

\ecoro

\bcoro

    Let $R$ be a PID.
    Then an $R$-module $M$ is injective iff it is {\emphcolor divisible}: for every nonzero $r\in R$ and $a\in M$ there is a $b\in M$ such that $a=br$.

\ecoro

\bproof

    If $M$ is injective then for $r\in R$ and $a\in M$ we define $f\colon R\to M$ by $f(1)=a$ and extend by linearity: $f(x)=xa$.
    And define $g\colon R\to R$ by $g(1)=r$ (so $g(x)=rx$).
    Since $R$ is an integral domain $g$ is injective, and thus we can lift $f$ to $f'\colon R\to M$ such that $f'g=f$.
    Then $f'g(1)=f'(r)=rf'(1)=a$ as desired.

    If $M$ is divisible then for every $x\in R$, a map $f\colon (x)\to M$ is determined by its value on $x$: $f(rx)=rf(x)$.
    We want to find a value $f(1)$ such that $xf(1)=f(x)$, which exists precisely because $M$ is divisible.
    And so we can extend $f$ to $f\colon R\to M$ by $f(r)=rf(1)$.
    All ideals of $R$ are of this form (principal) since $R$ is a PID, so by Baer $M$ is injective as desired.
    \qed

\eproof

In particular $\bQ,\bZ_{p^\infty}=\bZ[1/p]/\bZ,\bQ/\bZ$ are injective Abelian groups ($\bZ[1/p]$ are all rationals of the form $a/p^n$).

\bprop

    $\Ab$ has enough injectives.

\eprop

\bproof

    For an Abelian group $A$, define $I(A)$ to be the product of $\bQ/\bZ$ indexed by $\hom_\Ab(A,\bQ/\bZ)$.
    This is a product of injectives and is thus also injective.
    Define $e_A\colon A\to I(A)$ by $e_A(a)=(f(a))_{f\colon A\to\bQ/\bZ}$.

    We claim that $e_A$ is an injection, so we must show that $\ker e_A=0$.
    To do this we show that for every $0\neq a\in A$, there is some $f\colon A\to\bQ/\bZ$ such that $f(a)\neq0$.
    Let $\gen a\leq A$ be the subgroup generated by $a$, then by injectivity of $\bQ/\bZ$ we can lift any map $f\colon\gen a\to\bQ/\bZ$ to a map $A\to\bQ/\bZ$
    (following the inclusion $\gen a\subseteq A$).
    Now $\gen a\cong\bZ/m\bZ$ or $\bZ$ depending on $a$'s order; in the first case define $f(a)=1/m$ and in the second define $f(a)=1/2$ say.
    In both cases these are well-defined maps where $f(a)\neq0$, as desired.
    \qed

\eproof

Note that since $A\to(\bQ/\bZ)^{\hom(A,\bQ/\bZ)}$ is injective, if $\hom(A,\bQ/\bZ)=0$ the right-hand side is zero and so too is $A$.
Thus $A=0$ iff $\hom(A,\bQ/\bZ)=0$.

\bprop

    $I\in\A$ is injective iff every short exact sequence $0\to I\to A\to B\to0$ splits.

\eprop

\bproof

    By the duality between injectives and projectives.
    \qed

\eproof

\bdefn

    Let $M\in\A$, then a {\emphcolor right resolution} of $M$ is a cochain complex $I$ with $I^n=0$ for $n<0$ and a map $M\to I^0$ such that the augmented
    complex
    $$ 0 \to M \to I^0 \xtos d I^1 \xtos d \cdots $$
    is exact.
    This is the same as a quasi-isomorphism $M\to I$.
    This is a {\emphcolor injective resolution} if each $I_n$ is injective.

\edefn

A right resolution in $\A$ is the same as a left resolution in $A^\op$; an injective resolution in $\A$ is a projective resolution in $\A^\op$.
Thus by dualizing our previous results we have the following:

\blemm

    If $\A$ has enough injectives, every object in $\A$ has an injective resolution.

\elemm

\bthrm[title=Comparison theorem]

    Let $N\to I$ be an injective resolution of $N$, and $f'\colon M\to N$ a map in $\A$.
    Then for every resolution $M\to E$ there is a cochain map $f\colon E\to N$ lifting $f'$, unique up to chain homotopy equivalence.

\ethrm

\bprop

    $I\in\Ch(\A)$ is injective iff it is a split exact chain complex of injectives.
    Furthermore if $\A$ has enough injectives, so does $\Ch(\A)$.

\eprop

\bproof

    This follows immediately from previous results on projectives and the isomorphism $\Ch(\A)^\op\cong\Ch(\A^\op)$.
    This can be seen as follows: $\Ch(\A)$ can be viewed as the functor category $[\bZ_0,\A]$ (of additive functors) where $\bZ_0$ is the category whose objects are integers and $\hom(i,j)=\bZ$ when $j=i-1$
    and $0$ otherwise.
    Now in general $[\C,\D]^\op\cong[\C^\op,\D^\op]$ by mapping $F\colon\C\to\D$ to its opposite functor $F^\op\colon\C^\op\to\D^\op$ (which acts precisely as $F$ does).
    Note that natural transformations $F^\op\To G^\op$ are precisely natural transformations $G\To F$, so this is a well-defined functor and clearly an isomorphism.

    Note that $\bZ_0^\op\cong\bZ_0$ by mapping $n\mapsto-n$, and so
    $$ \Ch(\A)^\op \cong [\bZ_0,\A]^\op \cong [\bZ_0,\A^\op] \cong \Ch(\A^\op) $$
    as desired.

    Explicitly, we map a chain complex $C_\bul$ to its cochain complex $C^\bul$ (recall $C^n=C_{-n}$), and a chain map to its dual ($f^n=f_{-n}$).
    \qed

\eproof

We now will show that every category of modules has enough injectives.
Note that since the opposite of a category of modules is itself not a category of modules, this is not a trivial result.

If $A$ is an Abelian group and $B$ a left $R$-module, then $\hom_\Ab(B,A)$ can be given a right $R$-module structure by $fr\colon b\mapsto f(rb)$.

\blemm

    For every right $R$-module $M$, we have a natural isomorphism
    $$ \tau\colon\hom_\Ab(M,A) \to \hom_{\Mod_R}(M,\hom_\Ab(R,A)) $$
    given by $(\tau f)(m)\colon r\mapsto f(mr)$.

\elemm

\bproof

    We define $\tau$'s inverse $\mu$ as follows: for $g\colon M\to\hom_\Ab(R,A)$ define $(\mu g)(m)=g(m)(1)$.
    For $f\in\hom_\Ab(M,A)$ then $(\mu\tau f)(m)=(\tau f)(m)(1)=f(m)$ so $\mu\tau$ is the identity.
    And for $f\in\hom_{\Mod_R}(M,\hom_\Ab(R,A))$ then $(\tau\mu f)(m)\colon r\mapsto(\mu f)(mr)=f(mr)(1)$.
    Notice $f(m)(r)=(f(m)r)(1)=(fmr)(1)$ as desired, so $\tau\mu f=f$ as desired.
    \qed

\eproof

Thus $\hom_\Ab(R,\bul)\colon\Ab\to\Mod_R$ is {\emphcolor right-adjoint} to the forgetful functor $U\colon\Mod_R\to\Ab$:

\bdefn

    Functors $L\colon\A\tof\B\colon R$ form an {\emphcolor adjunction} if there is a natural isomorphism
    $$ \tau_{AB}\colon\hom_\B(L(A),B) \cong \hom_\A(A,R(B)) $$
    (i.e.\ a natural isomorphism of the bifunctors given by the compositions of the hom-functors and $L$ and $R$.)
    $L$ is the {\emphcolor left adjoint} and $R$ the {\emphcolor right adjoint}.
    This is commonly written $L\dashv R$.

\edefn

\bprop

    Suppose we have an adjunction of additive functors $L\colon\A\tof\B\colon R$, then
    \benum
        \item if $L$ is exact and $I\in\B$ injective then $R(I)\in\A$ is injective as well;
        \item dually if $R$ is exact and $P\in\A$ projective then $L(P)\in\B$ is projective.
    \eenum

\eprop

\bproof

    We must show that $\hom_\A(\bul,R(I))$ is exact, but by assumption this is isomorphic to $\hom_\B(L\bul,I)$.
    Since $I$ is injective, $\hom_\B(\bul,I)$ is exact, and the composition of exact functors is obviously exact.
    Duality follows because $L\dashv R$ iff $R^\op\dashv L^\op$.
    \qed

\eproof

We showed that $U\dashv\hom_\Ab(R,\bul)$ and the forgetful functor $U$ is clearly exact.
Thus $\hom_\Ab(R,I)$ is an injective $R$-module for any injective Abelian group $I$.

\bthrm

    $\Mod_R$ has enough injectives.

\ethrm

\bproof

    For an $R$-module $M$, define $I(M)$ to be the product of $I_0=\hom_\Ab(R,\bQ/\bZ)$ with itself indexed by $\hom_R(M,I_0)$.
    As explained, $I_0$ is injective and thus so is $I(M)$.
    Now define $e_M\colon M\to I(M)$ by $e_M(m)=(f(m))_{f\colon M\to I_0}$.
    Let $\alpha\colon R\to\bQ/\bZ$ be non-zero (which exists by extending a non-zero map $\gen1\to\bQ/\bZ$ where $\gen1$ is the subgroup generated by $1$, which exists as per
    our previous argument).
    Now for $m\in M$ define $f\colon mR\to I_0$ by $f(m)=\alpha$, and since $I_0$ is injective this extends to $M\to I_0$ where $f(m)\neq0$.
    Thus $e_M$ has trivial kernel and is thus an injection.
    \qed

\eproof

\bexam

    Consider $\sh(X)$, the category of sheaves (of Abelian groups) over a topological space $X$.
    We claim that $\sh(X)$ has enough injectives.

    For a sheaf $\F$ and a point $x\in X$, $\F$'s {\emphcolor stalk} at $x$ is the colimit $\F_x=\colim\set{\F(U)}[x\in U]$ (i.e.\ the limit of the restriction of $\F$ to the poset of open subsets containing $x$).
    This is functorial $\sh(X)\to\Ab$ for every $x\in X$ (since $\colim$ is).
    If $A$ is an Abelian group and $x\in X$, we define the {\emphcolor skyscraper sheaf} $x_*A$ at $x\in X$ to be the presheaf $(x_*A)(U)=A$ iff $x\in U$ and $0$ otherwise.
    $x_*$ acts on inclusions by either maps to zero or the identity, that is if $U\subseteq V$ then $\bul|_U\colon(x_*A)(V)\to(x_*A)(U)$ is the identity if $x\in U$ and the zero map otherwise.

    The skyscraper sheaf is indeed a sheaf.
    If $\set{U_i}$ is an open cover of $U$, suppose $x\in U$; then $a_i,a_j$ ($a_i\in(x_*A)(U_i)$) agreeing on their intersection means that $a_i=a_j\in A$ if $x\in U_i,U_j$.
    Let $a$ be the common element of all open sets containing $x$, this is the unique extension.
    Otherwise if $x\notin U$ all $a_i=0$, so $a=0$.

    Note that $\set{U}[x\in U]$ is a directed set and as such $\F_x$ is a direct limit which we can explicitly write out.
    $\F_x$ has underlying set $\bigsqcup\F(U)$ quotiented by an equivalence, where $s_i\in\F(U_i),s_j\in\F(U_j)$ are equated iff they have a common restriction (on some neighborhood of $x$).

    Now we claim that $\F_x\dashv(x_*A)$ (functorial in $x$) for every sheaf $\F$ and Abelian group $A$.
    That is, we need a natural isomorphism
    $$ \hom_\Ab(\F_x,A) \cong \hom_{\sh(X)}(\F,x_*A) $$
    Let $\phi\colon\F_x\to A$ be an Abelian morphism, then define $\phi'_U\colon\F(U)\to x_*(A)$ by $\phi'_U(s)=\phi([s])$ (where $[s]$ is the equivalence class of $s$ in $\F_x$), this defines a morphism of sheaves
    $\phi'\colon\F\to x_*A$.
    Indeed: if $x\in U\subseteq V$ then $\phi'_U(s|_U)=\phi([s])=\phi'_V(s)$ (the restriction of $x_*A(V)$ to $x_*A(U)$ is the identity.
    If $x\notin U$ then all maps compose to zero.

    Conversely given $\psi\colon\F\to x_*A$ defining $\hat\psi\colon\F_x\to A$ by $\hat\psi([s])=\psi_U(s)$ where $s\in\F(U)$.
    This is well-defined: if $s\in\F(U)$ and $s'\in\F(V)$ have the same restriction to a neighborhood of $x$, $W$ ($s|_W=s'|_W$) then because $\psi$ is a sheaf map we have that $\psi_U(s)|_W=\psi_V(s')|_W$.
    But restriction in the skyscraper sheaf is the identity: $\psi_U(s)=\psi_V(s')$ as desired.
    The composition of $\hat\bul$ and $\bul'$ gives the identities clearly, so we have a natural isomorphism, as desired.

    Note that $\F\mapsto\F_x$ is an exact functor: for a morphism of sheaves $\phi\colon\F\to\c G$ the induced map $\phi_x\colon\F_x\to\c G_x$ is given by $\phi_x([s])=[\phi_U(s)]$ for $s\in\F(U)$.
    The kernel of $\phi$ can be given by the sheaf $\ker\phi(U)=\ker\phi_U$, with restriction inherited from $\F$ (since $\ker\phi(U)\subseteq\F(U)$).
    And its image is given by $\im\phi(U)=\im\phi_U$, with restriction inherited from $\c G$.
    Then $\ker\phi_x=\set{[s]\in\F_x}[\hbox{$s|_U=0$ for some $x\in U$}]=\bigcup_{x\in U}[\ker\phi_U]$, and $\im\phi_x=\bigcup_{x\in U}[\im\phi_U]$.
    So if $\im\phi=\ker\psi$ we see that $\im\phi_x=\ker\psi_x$ as desired.

    Thus $x_*A$ is right-adjoint to an exact functor; it preserves injectives.
    So for every $x\in X$ choose an injective Abelian group, then $x_*A_x$ is injective in $\sh(X)$, and thus so is $\prod_{x\in X}x_*A_x$.

    Now given a sheaf $\F$, for every $x\in X$ choose an injection $\F_x\to I_x$ with $I_x$ injective.
    Composing $\F\to x_*\F_x$ ($\phi_U(s)=[s]$ if $x\in U$ and $0$ otherwise) with $x_*\F_x\to x_*I_x$ gives a map $\F\to\prod_{x\in X}x_*I_x=\c I$.
    This is an injection: for $0\neq s\in\F(U)$, note that there must be some $x\in V$ such that $s|_W\neq0$ for all $W\subseteq V$ (by the sheaf axiom), and so $\phi_V(s)\neq0$.
    Thus its kernel is zero.

    So we have shown $\sh(X)$ has enough injectives.

\eexam

\bexam

    Let $I$ be a small category and $\A$ Abelian and complete with enough injectives, we will show the functor category $\A^I=[I,\A]$ has enough injectives.

    Consider the {\emphcolor $k$th coordinate functor} for $k\in I$: map $A\colon I\to\A$ to $A(k)\in\A$.
    This is an exact functor.

    Now for $k\in I$ and $A\in\A$ define $k_*A\colon I\to A$ mapping $i\in I$ to
    $$ k_*A(i) = \prod_{\hom_I(i,k)}A $$
    and if $f\colon i\to j$ is a map in $I$, the map $k_*A(f)\colon k_*A(i)\to k_*A(j)$ is determined by $f^*\colon\hom(j,k)\to\hom(i,k)$.
    That is, the coordinate $k_*A(i)\to A$ of this map corresponding to $\phi\in\hom_I(i,k)$ is the projection of $k_*A(i)$ onto the factor corresponding to $f^*\phi=\phi f$.
    And for $f\colon A\to B\in\A$ then define $k_*f\colon k_*A\to k_*B$ which maps each component $A\in k_*A$ to the corresponding term $B\in k_*B$ via $f$.
    This makes $k_*$ a functor $\A\to\A^I$.

    Now, $k_*$ is right-adjoint to the $k$th coordinate functor:
    $$ \hom_\A(A(k), B) \cong \hom_{\A^I}(A,k_*B),\qquad A\in\A^I,B\in\A $$
    Map $\phi\colon A(k)\to B$ to $\phi'\colon A\To k_*B$ where $\phi'_i\colon A(i)\to k_*B(i)$ is defined by mapping $A(i)\xto{A(f)}A(k)\xto\phi B$ for each component $f\colon i\to k$.
    And for $\psi\colon A\To k_*B$ map to $\hat\psi\colon A(k)\to B$ by the composition $A(k)\xto{\psi_k}k_*B(k)\xto{\pi_\id}B$ where $\pi_\id$ is the projection to the component associated with $\id\in\hom(k,k)$.

    Thus $k_*$ preserves injectives.
    For $F\in\A^I$, for each $k\in I$ embed $F(k)\to A_k$ for $A_k$ injective.
    Then $E=\prod_{k\in I}k_*A_k$ exists (functor categories over complete categories are complete) and is injective.
    Further, $\eta\colon F\To E$ is given by $\eta_i\colon F(i)\to\prod_{k\in I}k_*A_k(i)$ whose projection $F(i)\to A_k$ for $\phi\colon i\to k$ is given by the composition $F(i)\xto{F\phi}F(k)\to A_k$.
    This is a monomorphism.

    Thus when $\A$ is complete Abelian and has enough injectives, so does $\A^I$.

    Conversely if $\A$ is cocomplete and has enough projectives, $\A^\op$ has enough injectives and thus so does $[I^\op,\A^\op]\cong(\A^I)^\op$.
    Thus $\A^I$ has enough projectives.

\eexam

\subsection{Left derived functors}

\bdefn

    Let $F\colon\A\to\B$ be a right exact functor between Abelian categories, where $\A$ has enough projectives.
    Define the {\emphcolor left derived functors} of $F$, $\set{L_iF}_{i\geq0}$ as follows.
    If $A$ is an object of $\A$, choose a projective resolution $P\to A$ and define $L_iF(A)=H_i(F(P))$.
    For a map $f\colon A\to B$ in $\A$, suppose we chose a resolution $Q\to B$ then by comparison we have a unique extension of $f$ to a chain map between the complexes (up to chain homotopy),
    then $Ff$ is a chain map and as such we have a unique map $L_iF(f)$ from $H_i(F(P))\to H_i(F(Q))$ (since chain homotopic maps induce the same morphisms on homology modules).

\edefn

Note that $H_i(F(P))$ corresponds to the homology module of the complex $\cdots\to F(P_1)\to F(P_0)\to0$, not the augmented complex containing $A$.

These are indeed functors: clearly $\id\colon A\to A$ lifts to $\id\colon P\to P$ and thus $L_iF(\id)$ is the identity.
And given $P\to A,P'\to A',P''\to A''$ and maps $A\xtos fA'\xtos{f'}A''$, $f$ lifts to $P\to P'$ and $f'$ lifts to $P'\to P''$, then the composition of these lifts is a lift of $f'\circ f$.
They are additive functors, too: for $f,g\colon A\to B$ if $\tilde f,\tilde g$ are their lifts then $\tilde f+\tilde g$ lifts $f+g$.

Now notice that since $F$ is right exact, $F(P_1)\xto dF(P_0)\xto{d'}F(A)\to0$ is exact, and as such such $H_0(F(P_0))=F(P_0)/\im d$ and we know $\im d=\ker d'$.
Since $d'$ is a projection we have $H_0(F(P_0))=F(P_0)/\ker d'\cong F(A)$.
Thus $L_0F(A)\cong F(A)$, and since $L_0F(f)$ is just the induced map of $f$, we have that $L_0\cong F$ naturally.

\blemm

    Let $F$ have left derived functors $L_iF$, then $L_0F\cong F$.

\elemm

Also we show the naturality of $L_iF$.

\blemm

    Different choices of projective resolutions result in naturally isomorphic left derived functors.
    That is, left derived functors are well-defined up to a unique natural isomorphism.

\elemm

\bproof

    Let $P\to A$ and $Q\to A$ be two projective resolutions.
    Then there are chain maps $f\colon P\tof Q\colon g$ lifting $\id_A$, and these are unique up to chain homotopy.
    But then $gf\colon P\to P$ and $fg\colon Q\to Q$ are lifts of $\id_A$, and since $\id_P,\id_Q$ are lifts as well, $gf$ and $fg$ are chain homotopic to the identities.
    Consider then the induced maps $F(f)_*\colon H_i(F(P))\tof H_i(F(Q))\colon F(g)_*$.
    Additive functors preserve chain homotopy, so $F(fg),F(gf)$ are chain homotopic to the identities, and as such $F(fg)_*=F(f)_*F(g)_*$ is the identity, similar for $gf$.
    Thus we have an isomorphism $H_iF(P)\cong H_iF(Q)$ as desired.
    \qed

\eproof

\bcoro

    If $A$ is projective, $L_iF(A)=0$ for $i\neq0$.

\ecoro

\bproof

    There is a simple projective resolution of $A$: $\cdots\to0\to0\to A\xtos\id A\to0$.
    Then applying $F$ gives us the chain complex $\cdots\to0\to0\to FA\to0$, whose nonzero homology groups are all $0$.
    \qed

\eproof

\bdefn

    An object $Q\in\A$ is {\emphcolor $F$-acyclic} if $L_iFQ=0$ for all $i\neq0$.
    An {\emphcolor $F$-acyclic resolution} of $A\in\A$ is a left resolution $Q\to A$ where the $Q_i$s are all $F$-acyclic.

\edefn

We will later show that $L_iF(A)\cong H_i(F(Q))$ for any $F$-acyclic resolution $Q$ of $A$.

\bthrm

    The derived functors $L_\bul F$ form a homological $\delta$-functor.

\ethrm

\bproof

    Given a short exact sequence $0\to A'\to A\to A''\to0$ choose projective resolutions $P'\to A'$ and $P''\to A''$.
    By the horseshoe lemma there is a projective resolution $P\to A$ fitting into a short exact sequence $0\to P'\to P\to P''\to0$ of projective complexes in $\A$.
    Since $P''_n$ is projective, each $0\to P'_n\to P_n\to P''_n\to0$ is split exact, and thus so is $0\to FP'_n\to FP_n\to FP''_n\to0$, in particular we have a short exact sequence
    $0\to FP'\to FP\to FP''\to0$.
    We thus have a long exact homology sequence
    $$ \cdots\xtos\partial L_iF(A') \to L_iF(A) \to L_iF(A'') \xtos\partial L_{i-1}F(A') \to L_{i-1}F(A) \to L_{i-1}F(A'') \xtos\partial\cdots $$

    These $\partial$s are natural (and thus form the $\delta$-morphisms of the $\delta$-functor).
    Indeed, suppose we have a morphism of exact sequences

    \centerline{\drawdiagram{
        $0$&$A'$&$A$&$A''$&$0$\cr
        $0$&$B'$&$B$&$B''$&$0$\cr
    }{
        \diagarrow{from={1,1}, to={1,2}}
        \diagarrow{from={1,2}, to={1,3}}
        \diagarrow{from={1,3}, to={1,4}}
        \diagarrow{from={1,4}, to={1,5}}
        \diagarrow{from={2,1}, to={2,2}}
        \diagarrow{from={2,2}, to={2,3}, text=$i_B$, y distance=.25cm}
        \diagarrow{from={2,3}, to={2,4}, text=$\pi_B$, y distance=.25cm}
        \diagarrow{from={2,4}, to={2,5}}
        \diagarrow{from={1,2}, to={2,2}, text=$f'$, x distance=-.25cm}
        \diagarrow{from={1,3}, to={2,3}, text=$f$, x distance=-.25cm}
        \diagarrow{from={1,4}, to={2,4}, text=$f''$, x distance=-.25cm}
    }}

    and choose projective resolutions $\epsilon'\colon P'\to A',\epsilon\colon P''\to A'',\eta'\colon Q'\to B',\eta''\colon Q''\to B''$.
    The connecting morphisms are obtained via the horseshoe lemma, which give $\epsilon\colon P\to A,\eta\colon Q\to B$.
    The horseshoe lemma gives us $\epsilon$ is the direct sum of the lift of $\epsilon''$ with $i_A\epsilon'$.
    The comparison lemma gives us unique lifts of $f',f''$ to $F'\colon P'\to Q'$ and $F''\colon P''\to Q''$.
    We show that there is a lift of $f$ to $F\colon P\to Q$ such that the following diagram commutes

    \centerline{\drawdiagram{
        $0$&$P'$&$P$&$P''$&$0$\cr
        $0$&$Q'$&$Q$&$Q''$&$0$\cr
    }{
        \diagarrow{from={1,1}, to={1,2}}
        \diagarrow{from={1,2}, to={1,3}}
        \diagarrow{from={1,3}, to={1,4}}
        \diagarrow{from={1,4}, to={1,5}}
        \diagarrow{from={2,1}, to={2,2}}
        \diagarrow{from={2,2}, to={2,3}}
        \diagarrow{from={2,3}, to={2,4}}
        \diagarrow{from={2,4}, to={2,5}}
        \diagarrow{from={1,2}, to={2,2}, text=$F'$, x distance=-.25cm}
        \diagarrow{from={1,3}, to={2,3}, text=$F$, x distance=-.25cm}
        \diagarrow{from={1,4}, to={2,4}, text=$F''$, x distance=-.25cm}
    }}

    To construct $F$, we define $\gamma_n\colon P''_n\to Q'_n$ (not a chain map) such that
    $$ F_n = \pmatrix{F'_n&\gamma_n\cr0&F''_n} \colon P'_n \oplus P''_n \to Q'_n \oplus Q''_n,\qquad
    F(p',p'') = (F'p'+\gamma(p''),F''(p'')) $$
    If $F$ is a chain map, then this diagram will indeed commute.
    In order for $F$ to lift $f$ we must have that $\eta F_0-f\epsilon$ is zero on $P_0=P_0'\oplus P_0''$.
    This is true for any $\gamma$ on $P_0'$, so we must require just that
    $$ i_B\eta'\gamma_0 = f\lambda_P - \lambda_QF''_0 $$
    where $\lambda_P,\lambda_Q$ are the restrictions of $\epsilon,\eta$ to $P_0''$ and $Q_0''$, $i_B$ is the inclusion of $B'\to B$.
    Note that in $B''$ we have
    $$ \pi_B(f\lambda-\lambda F_0'') = f''\pi_A\lambda_P - \pi_B\lambda F_0'' = f''\epsilon'' - \eta''F_0'' = 0 $$
    thus there is a map $\beta\colon P_0''\to B'$ such that $i_B\beta=f\lambda-\lambda F_0''$ (since it is zero on $B''$).
    Lift $\beta$ to $\gamma_0\colon P_0''\to Q_0'$: $\beta=\eta'\gamma_0$.
    Now we have
    $$ dF - Fd = \pmatrix{d'F'-F'd'& (d'\gamma-\gamma d''+\lambda F''-F'\lambda')\cr0&d''F''-F''d''} $$
    and for $F$ to be a chain map, this must vanish.
    So we must have $d'\gamma=\gamma d''-\lambda F''+F'\lambda'$.
    This can be constructed inductively.
    \qed

\eproof

\bprop[title=Dimension shifting]

    If $0\to M\to P\to A\to0$ is exact with $P$ projective (or $F$-acyclic), then $L_iF(A)\cong L_{i-1}F(M)$ for $i\geq2$,
    and $L_1F(A)$ is the kernel of $F(M)\to F(P)$.

    In general if we have an exact sequence
    $$ 0 \to M_m \to P_m \to P_{m-1} \to \cdots \to P_0 \to A \to 0 $$
    with $P_i$ projective (or in general $F$-acyclic), then $L_iF(A)\cong L_{i-m-1}F(M_m)$ for $i\geq m+2$, and $L_{m+1}F(A)$ is the kernel of
    $F(M_m)\to F(P_m)$.

\eprop

\bproof

    If $0\to M\to P\to A\to0$ is short exact, then since derived functors are $\delta$-functors, we have a long exact sequence
    $$ \cdots\to L_iF(P) \to L_iF(A) \to L_{i-1}F(M) \to L_{i-1}F(P) \to \cdots $$
    and for $i\geq2$ we have that $L_iF(P)=L_{i-1}F(P)=0$ and so we have the isomorphism.
    For $i=1$ (recalling $L_0F\cong F$) this gives us $0\to L_1F(A)\to F(M)\to F(P)\to0$ as desired.

    Now we induct on $m$, we have shown it for $m=0$.
    Given an exact sequence $0\to M_m\to P_m\to\cdots\to P_0\to A\to0$, consider the two exact sequences
    $$ 0 \to P_m/M_m \to P_{m-1} \to \cdots P_0\to A\to0,\qquad 0\to M_m \to P_m \to P_m/M_m\to 0 $$
    then inductively we have $L_iF(A)\cong L_{i-m}F(P_m/M_m)\cong L_{i-m-1}F(M_m)$ (the second isomorphism due to our base case).
    \qed

\eproof

In particular suppose we have an $F$-acyclic resolution
$$ \cdots P_{m+1} \xto{f_{m+1}} P_m \xto{f_m} \cdots \to P_0 \to A \to 0 $$
we can turn this into
$$ 0 \to \coker(f_{m+1}) \to P_{m-1} \cdots \to P_0 \to A \to 0 $$
so that $L_mF(A)$ is the kernel of the map $F\coker(f_{m+1})\to FP_{m-1}$, by our above proposition.
Since $F$ is right-exact, it preserves cokernels, and in general the homology group of $A\to B\to C$ is the kernel of $\coker(A\to B)\to C$, as desired.

\bprop

    If $P\to A$ is an $F$-acyclic resolution of $A$, then $L_mF(A)\cong H_m(FP)$.

\eprop

\bthrm

    If $\A$ has enough projectives, the left derived functors of any right exact functor $F\colon\A\to\B$ form a universal homological $\delta$-functor.

\ethrm

\bproof

    Suppose we have another homological $\delta$-functor $T_*$ and a natural map $\phi_0\colon T_0\to F\cong L_0F$.
    We need to show that $\phi_0$ can be uniquely extended to a map $\phi_\bul\colon T_\bul\to L_\bul F$.
    We construct $\phi_n$ inductively: suppose $\phi_i$ is constructed for $0\leq i<n$, such that they commute with the appropriate $\delta$s.

    Given $A\in\A$, choose an exact sequence $0\to K\to P\to A\to0$ with $P$ projective.
    Since $L_nF(P)=0$ we have a commutative diagram with exact rows:

    \centerline{\drawdiagram{
        &$T_nA$&$T_{n-1}K$&$T_{n-1}P$\cr
        $0$&$L_nFA$&$L_{n-1}FK$&$L_{n-1}FP$\cr
    }{
        \diagarrow{from={1,2}, to={1,3}, text=$\delta_n$, y distance=-.25cm}
        \diagarrow{from={2,2}, to={2,3}, text=$\delta_n$, y distance=.25cm}
        \diagarrow{from={1,3}, to={2,3}, text=$\phi_{n-1}$, x distance=.25cm}
        \diagarrow{from={1,4}, to={2,4}, text=$\phi_{n-1}$, x distance=.25cm}
    }}

    Since the bottom $\delta_n$ is monic, $\delta_n\circ\phi_n=\phi_{n-1}\circ\delta_n$ has a unique solution (explicitly, if a solution exists it is unique,
    we can check that such a solution exists in $\Mod_R$ via diagram chase).
    So we have defined $\phi_n$ on $A$: we must now show that the resulting map is natural $T_n\to L_nF$.
    Note that our $\phi_n$ ostensibly depends on choice of $K,P$; we will show that this is not the case.

    So suppose we have a map $f\colon A'\to A$ and an exact sequence $0\to K'\to P'\to A'\to0$ with $P'$ projective.
    Since $P'$ is projective, we can lift $P'\to A$ (composing $P'\to A'$ with $f$) to $g\colon P'\to P$ (via the epic $P\to A$).
    This in turn induces a map $h\colon K'\to K$ such that the diagram commutes

    \centerline{\drawdiagram{
        $0$&$K'$&$P'$&$A'$&$0$\cr
        $0$&$K$&$P$&$A$&$0$
    }{
        \diagarrow{from={1,1}, to={1,2}}
        \diagarrow{from={1,2}, to={1,3}}
        \diagarrow{from={1,3}, to={1,4}}
        \diagarrow{from={1,4}, to={1,5}}
        \diagarrow{from={2,1}, to={2,2}}
        \diagarrow{from={2,2}, to={2,3}}
        \diagarrow{from={2,3}, to={2,4}}
        \diagarrow{from={2,4}, to={2,5}}
        \diagarrow{from={1,2}, to={2,2}, text=$h$, x distance=.25cm}
        \diagarrow{from={1,3}, to={2,3}, text=$g$, x distance=.25cm}
        \diagarrow{from={1,4}, to={2,4}, text=$f$, x distance=.25cm}
    }}

    We now show that $\phi_n$ commutes with $f$.
    Note that in the following diagram, each small quadrilateral commutes:

    \centerline{\drawdiagram{
        $T_n(A')$&&&$T_n(A)$\cr
        &$T_{n-1}(K')$&$T_{n-1}(K)$\cr
        &$L_{n-1}F(K')$&$L_{n-1}F(K)$\cr
        $L_nF(A')$&&&$L_nF(A)$\cr
    }{
        \diagarrow{from={1,1}, to={1,4}, text=$T_nf$, y distance=-.25cm}
        \diagarrow{from={4,1}, to={4,4}, text=$L_nFf$, y distance=.25cm}
        \diagarrow{from={2,2}, to={2,3}, text=$T_nh$, y distance=-.25cm}
        \diagarrow{from={3,2}, to={3,3}, text=$L_nFh$, y distance=.25cm}
        \diagarrow{from={1,1}, to={4,1}, text=$\phi_n(A')$, x distance=-.25cm}
        \diagarrow{from={1,4}, to={4,4}, text=$\phi_n(A)$, x distance=.25cm}
        \diagarrow{from={2,2}, to={3,2}, text=$\phi_{n-1}$, x distance=-.25cm}
        \diagarrow{from={2,3}, to={3,3}, text=$\phi_{n-1}$, x distance=.25cm}
        \diagarrow{from={1,1}, to={2,2}, text=$\delta$, x distance=.25cm}
        \diagarrow{from={1,4}, to={2,3}, text=$\delta$, x distance=-.25cm}
        \diagarrow{from={4,1}, to={3,2}, text=$\delta$, x distance=-.25cm}
        \diagarrow{from={4,4}, to={3,3}, text=$\delta$, x distance=.25cm}
    }}

    And so we have
    $$ \delta\circ L_n(f)\circ\phi_n(A') = \delta\circ\phi_n(A)\circ T_n(f) $$
    both of these $\delta$s are $L_nF(A)\to L_{n-1}F(K)$ which is monic, and thus we can cancel it from both sides.
    Thus we have that $\phi_n$ is indeed natural.
    Note that choosing $A=A'$ and $f=\id_A$ we see that $\phi_n(A)$ is independent on choice of short exact sequence.

    We also need to show that $\phi_n$ commutes with $\delta_n$.
    Given a short exact sequence $0\to A'\to A\to A''\to0$, choose a short exact sequence $0\to K''\to P''\to A''\to0$ with $P''$ projective.
    Then we can construct $f,g$ making the following diagram commute

    \centerline{\drawdiagram{
        $0$&$K''$&$P''$&$A''$&$0$\cr
        $0$&$A'$&$A$&$A''$&$0$
    }{
        \diagarrow{from={1,1}, to={1,2}}
        \diagarrow{from={1,2}, to={1,3}}
        \diagarrow{from={1,3}, to={1,4}}
        \diagarrow{from={1,4}, to={1,5}}
        \diagarrow{from={2,1}, to={2,2}}
        \diagarrow{from={2,2}, to={2,3}}
        \diagarrow{from={2,3}, to={2,4}}
        \diagarrow{from={2,4}, to={2,5}}
        \diagarrow{from={1,2}, to={2,2}, text=$g$, x distance=.25cm}
        \diagarrow{from={1,3}, to={2,3}, text=$f$, x distance=.25cm}
        \diagarrow{from={1,4}, to={2,4}, text={$=$}, x distance=.25cm}
    }}

    This gives us a commutative diagram

    \centerline{\drawdiagram{
        $T_n(A'')$&$T_{n-1}(K'')$&$T_{n-1}(A')$\cr
        $L_nF(A'')$&$L_{n-1}F(K'')$&$L_{n-1}F(A')$\cr
    }{
        \diagarrow{from={1,1}, to={1,2}, text=$\delta$, y distance=-.25cm}
        \diagarrow{from={1,2}, to={1,3}, text=$Tg$, y distance=-.25cm}
        \diagarrow{from={2,1}, to={2,2}, text=$\delta$, y distance=.25cm}
        \diagarrow{from={2,2}, to={2,3}, text=$LF(g)$, y distance=.25cm}
        \diagarrow{from={1,1}, to={2,1}, text=$\phi_n$, x distance=.25cm}
        \diagarrow{from={1,2}, to={2,2}, text=$\phi_{n-1}$, x distance=.25cm}
        \diagarrow{from={1,3}, to={2,3}, text=$\phi_{n-1}$, x distance=.25cm}
    }}

    the horizontal composites are the $\delta$s of $0\to A'\to A\to A''\to0$ (by the $\delta$-functorality of $T_\bul$, $L_\bul F$), giving us the desired commutativity of $\phi_n$ with $\delta$.
    \qed

\eproof

\bdefn

    An additive functor $F\colon\A\to\B$ is {\emphcolor effaceable} if for every $A\in\A$ there is a monomorphism $u\colon A\to I$ such that $F(u)=0$.
    $F$ is {\emphcolor coeffaceable} if for every $A$ there is an epimorphism $u\colon P\to A$ such that $F(u)=0$.

    A $\delta$-functor is (co)effaceable if every component (other than $0$) is (co)effaceable.

\edefn

\bthrm

    A homological coeffaceable $\delta$-functor is universal.
    Dually a cohomological effaceable $\delta$-functor is universal.

\ethrm

\bproof

    In the above proof, instead of choosing $P$ projective choose $P$ with $P\to A$ epic whose image is $0$.
    The proof then proceeds the same.
    \qed

\eproof

Note that left derived functor are coeffaceable (choosing $P\to A$ with $P$ projective), so this is a generalization of the above theorem.

Furthermore, homology $H_\bul\colon\Ch_{\geq0}(\A)\to\A$ is coeffaceable.
Indeed, recall that a map $f\colon B\to C$ induces a short exact sequence $0\to C\to\cone(f)\to B[-1]\to0$.
In particular for a chain complex $A$ we have an epic $\sigma_{\geq0}\cone(A)[1]\to A$ which maps $(a,a')\mapsto-a$.
Since $\id_A$ is a quasi-isomorphism we have that $\cone(A)$ is acyclic and in particular the induced homology map of this epimorphism is zero (except at the zero index, which is unimportant).
So we have shown

\bthrm

    Homology and cohomology are universal $\delta$-functors $\Ch_{\geq0}(\A)\to\A$ and $\Ch^{\geq0}(\A)\to\A$ respectively.

\ethrm

\subsection{Right derived functors}

\bdefn

    Let $F\colon\A\to\B$ be a left exact functor between additive categories.
    If $\A$ has enough injectives, we define $F$'s {\emphcolor right derived functors} $R^iF\colon\A\to\B$ as follows.
    For $A\in\A$, choose an injective resolution $A\to I$ and define
    $$ R^iF(A) = H^i(F(I)) $$

\edefn

Note that since $0\to F(A)\to F(I^0)\to F(I^1)$ is exact, we have that $R^0F\cong F$.

Furthermore, $F\colon\A\to\B$ defines a functor $F^\op\colon\A^\op\to\B^\op$ and $\A^\op$ has enough projectives so we have left derived functors $L_iF^\op$.
Since $I$ is a projective resolution of $A$ in $\A^\op$ we clearly have that
$$ R^iF = (L_iF^\op)^\op $$
thus all results about left derived functors hold dually for right derived functors.
In particular $R^iF$'s construction is independent on choice of injective resolution, $R^\bul F$ is a universal cohomological $\delta$-functor, and $R^iF(I)=0$ for injective $I$ when $i\neq0$.

\bdefn

    For an $R$-module $A$, the functor $\hom_R(A,\bul)$ is left-exact.
    Thus it has right derived functors, called {\emphcolor Ext groups}:
    $$ \ext^i_R(A,\bul) = R^i\hom_R(A,\bul) $$

\edefn

In particular we have $\ext^0_R(A,B)=\hom(A,B)$.

\bprop

    The following are equivalent:
    \benum
        \item $B$ is injective,
        \item $\hom_R(\bul,B)$ is exact,
        \item $\ext^i_R(A,B)$ vanishes for all $i\neq0$ and all $A$,
        \item $\ext^1_R(A,B)$ vanishes for all $A$.
    \eenum

\eprop

\bproof

    As already explained $(1)$ and $(2)$ are equivalent, and these imply the rest.
    Now if $\ext^1_R(\bul,B)=0$ then consider a short exact sequence $0\to C\to D\to E\to0$ then we have a long exact sequence
    $$ 0 \to \hom(C,B) \to \hom(D,B) \to \hom(E,B) \to \ext^1(C,B) = 0 $$
    thus $\hom(\bul,B)$ is exact as desired.
    \qed

\eproof

Dually we have

\bprop

    The following are equivalent:
    \benum
        \item $A$ is projective,
        \item $\hom_R(A,\bul)$ is exact,
        \item $\ext^i_R(A,B)$ vanishes for all $i\neq0$ and all $B$,
        \item $\ext^1(A,B)$ vanishes for all $B$.
    \eenum

\eprop

We can consider derived functors of contravariant functors.
For example, left $F\colon\A^\op\to\B$ be left exact, then if $\A$ has enough projectives ($\A^\op$ has enough injectives), $F$ has right derived functors via cohomology of
$F(P)$ for $P\to A$ a projective resolution.

In particular we have left exact $\hom_R(\bul,B)$ and thus right derived functors $R^i\hom_R(\bul,B)$.
We will show later that these are not new:
$$ R^\bul\hom_R(\bul,B)(A) \cong R^\bul\hom_R(A,\bul)(B) = \ext_R^\bul(A,B) $$

\bexam

    Consider the {\emphcolor global sections} functor $\Gamma\colon\sh(X)\to\Ab$ via $\Gamma(\F)=\F(X)$.
    It can be shown that this is right adjoint to the constant sheaves functor, and in particular is left exact.
    Thus $\Gamma$ has right derived functors, called the {\emphcolor cohomology functors} on $X$:
    $$ H^i(X,\F) = R^i\Gamma(\F) $$

\eexam

\subsection{Adjointness and exactness}

Note that we can characterize (left/right) exactness purely categorically:

\blemm

    A functor $F\colon\A\to\B$ is left exact iff it preserves finite limits, right exact iff it preserves finite colimits, and exact iff it preserves finite
    limits and colimits.

\elemm

\bproof

    It is a classic result that a functor is additive iff it preserves finite biproducts, and note that an exact diagram $0\to A\to B\to C$ is simply a kernel diagram.
    So preserving diagrams of this form is the same as preserving kernels, and thus a functor is left exact iff it preserves finite limits (finite biproducts and kernels).
    Similarly exact $A\to B\to C\to0$ are cokernel diagrams.
    \qed

\eproof

This result allows us to generalize the concept of exactness to arbitrary categories, but this is not necessary here.
More importantly (for us) it gives us the following.

\bcoro

    Left adjoints are right exact, right adjoints are left exact.

\ecoro

\bproof

    Left adjoints preserve (arbitrary) colimits, and right adjoints preserve (arbitrary) limits.
    \qed

\eproof

In particular (assuming the existence of enough projectives/injectives) left adjoints have left derived functors and right adjoints have right derived functors.

Another interesting result is the following:

\blemm

    The following are equivalent for an exact functor $F\colon\A\to\B$:
    \benum
        \item $F$ reflects exact sequences: if $FA\to FB\to FC$ is exact, so is $A\to B\to C$.
        \item $F$ is faithful.
        \item $F$ reflects zero morphisms: if $Ff=0$ then $f=0$.
        \item $F$ reflects zero objects: if $FA=0$ then $A=0$.
    \eenum

\elemm

\bproof

    Clearly $F$ reflects zero morphisms iff it is faithful as $Ff=Fg$ iff $F(f-g)=0$ iff $f=g$.
    If $F$ reflects zero morphisms then suppose $FA=0$; then $F\id_A=0$ and so $\id_A=0$, meaning $A=0$.

    If $F$ reflects zero objects, suppose $f\colon A\to B$ is non-zero: we want to show $Ff$ is non-zero as well.
    Consider the exact sequence
    $$ 0 \to \ker f \to A \xtos f B \to \coker f \to 0 $$
    which maps to the exact sequence
    $$ 0 \to F\ker f \to FA \xtos{Ff} FB \to F\coker f \to 0 $$
    $0$ is the unique map whose kernel and cokernel are isomorphisms, so we want to show that the induced maps $F\ker f\to FA,FB\to F\coker f$ are not isomorphisms.
    But indeed, $\coker(F\ker f\to FA)\cong F\coker(\ker f\to A)$ is non-zero since $\ker f\to A$ is non-zero; so $F\ker f\to FA$ is not an isomorphism.
    Similarly $\ker(FB\to F\coker f)\cong F\ker(B\to\coker f)$ is non-zero and thus $FB\to F\coker f$ is not an isomorphism.
    So $Ff$ is non-zero as well, meaning $F$ reflects zero morphisms if it reflects zero objects.

    Thus we have shown the equivalence of $(2),(3),(4)$.
    If $F$ reflects exact sequences then it reflects zero objects, as if $FA=0$ then $0\to FA\to0$ is exact and so $0\to A\to0$ is as well.

    And if $F$ is faithful then it must reflect exact sequences: it is sufficient to show that $F$ reflects kernels and cokernels.
    Let $A\xto fB\xto gC$ such that $Ff=\ker Fg$, and let $h\colon K\to B$ be a kernel of $g$.
    Since $F(gf)=0$ and $F$ is faithful, $gf=0$, so there is a unique $i\colon K\to A$ such that $f=hi$; we want to show $i$ is an isomorphism.
    Similarly since $F(gh)=0$ there is a unique $j\colon FA\to FK$ such that $Fh=(Ff)j$.
    Then we have that
    $$ F(f)jF(i) = F(h)F(i) = F(f) $$
    and since $F(f)$ is monic we have $jF(i)=\id$, in particular $F(i)$ is monic.
    Faithfulness reflects monomorphisms, so $i$ is monic as well.

    Let $p\colon K\to K/A$ be a cokernel of $i$, and $q\colon B\to B/A$ a cokernel of $f$.
    Now $qhi=qf=0$, which means there exists a unique $h'\colon K/A\to B/A$ such that $qh=h'p$.
    Since $h$ is monic, so is $h'$.
    And note
    $$ F(h'p) = F(qh) = F(q)F(f)j = 0 $$
    So $h'p=0$ and thus $p=0$ since $h'$ is monic.
    But $p=\coker i$, thus $i$ is an epic monic, thus an isomorphism as desired.
    \qed

\eproof

Recall that in particular we have the tensor-hom adjunction $\bul\otimes_RB\colon\Mod_R\tof\Ab\colon\hom_\Ab(B,\bul)$ for all $B\in{}_R\Mod$.
If $B$ is an $(R,S)$-bimodule then this adjunction lands in $\Mod_S$ and not just $\Ab$.
So $\bul\otimes_RB$ is left adjoint and thus right exact.

\bdefn

    For an $(R,S)$-bimodule $B$, define its {\emphcolor Tor modules} to be the right $S$-modules
    $$ \tor^R_n(A,B) = (L_n(\bul\otimes_RB))(A) $$
    for $A\in\Mod_R$.

\edefn

If $S'\subseteq S$ is a subring then for an $(R,S)$-bimodule $B$ we can also view it as an $(R,S')$-bimodule.
Thus $\tor^R_n(\bul,B)$ can be viewed to land in either $\Mod_S$ or $\Mod_{S'}$.
But the forgetful functor $U\colon\Mod_S\to\Mod_{S'}$ is exact, and thus
$$ U(L_\bul(\otimes_RB)) \cong L_\bul(U\otimes_RB) $$
thus the underlying $S'$-module is the same in either case.

Note that for a right $R$-module $A$, $A\otimes_R\bul$ is also right-exact (due to a similar adjunction), and thus we can construct its left derived functors
as well.
We will show soon that $L_\bul(A\otimes_R\bul)(B)\cong L_\bul(\bul\otimes_RB)(A)$, so this yields nothing new.

\bexam

    Consider the inclusion functor $I\colon\sh(X)\to\Psh(X)$.
    It has a left adjoint called {\emphcolor sheafification}, and is thus left exact and has right derived functors.
    Further consider the exact ``section'' functors $\Psh(X)\to\Ab$ mapping a presheaf $\F$ to $\F(U)$, call this $\Gamma_U$.
    Then we have
    $$ \Gamma_UR^n(I) \cong R^n(\Gamma_U\circ I) $$
    note that $\Gamma_U$ is just the section functor restricted to sheaves.
    $\Gamma_U\circ I$ just maps a sheaf $\F$ to $\Gamma\F|_U$ (the global section of the restriction of $\F$ to $U$), and thus by definition
    $$ \Gamma_UR^n(I) \cong H^n(U,\F|_U) $$
    meaning that $R^n(I)$ maps a sheaf $\F$ to the presheaf $U\mapsto H^n(U,\F|_U)$.

\eexam

\bexam

    Let $f\colon X\to Y$ be continuous, then define the {\emphcolor direct image} functor $f_*\colon\Psh(X)\to\Psh(Y)$ by $(f_*\F)(V)=\F(f^{-1}V)$
    for $V$ open in $Y$.
    For $V'\subseteq V$, we have $f^{-1}V'\subseteq f^{-1}V$ so this defines a presheaf: map $s\in f_*\F(V)$ to $s|_{f^{-1}V'}\in f_*\F(V')$.
    Note that if $\F\in\sh(X)$ is a sheaf, then consider the exact diagram for an open cover $\set{V_i}$ of $V\subseteq Y$.
    $\set{f^{-1}V_i}$ is an open cover of $f^{-1}V$ and so we have an exact diagram
    $$ 0 \to \F(f^{-1}V) \to \prod_i\F(f^{-1}V_i) \to \prod_{i,j}\F(f^{-1}V_i\cap f^{-1}V_j) $$
    The maps are precisely the maps in the diagram for $f_*\F$, thus $f_*\F$ is a sheaf.
    So $f_*$ restricts to a map $\sh(X)\to\sh(Y)$.

    We also define the {\emphcolor inverse image} functor $f^{-1}\colon\sh(Y)\to\sh(X)$ to map $\G\in\sh(Y)$ to the sheafification of the presheaf
    mapping $U$ to the direct colimit $\colim\G(V)$ over the poset of all $V\subseteq Y$ containing $f(U)$.
    Then we have that $f^{-1}\dashv f_*$ form an adjunction.

    Now, let $*$ be the one-point topological space.
    Then we have $\sh(*)\cong\Ab$.
    Consider the unique map $f\colon X\to*$.
    The direct image functor $f_*\colon\sh(X)\to\Ab$ maps $\F\in\sh(X)$ to $\F(f^{-1}*)=\F(X)$, i.e.\ $f_*$ is the global section functor.
    The inverse image functor maps $A\in\Ab$ to the sheafification of $\F(U)=A$, i.e.\ the constant sheaf.
    Thus we have that ${\rm const}\dashv\Gamma$: the constant sheaf functor is a left adjoint to the global section functor.
    This shows us that $\Gamma$ is indeed left exact and so sheaf cohomology is well-defined.

    For $x\in X$ consider the map $x\colon*\to X$ mapping the point to $x\in X$.
    Then $x_*\colon\Ab\to\sh(X)$ maps an Abelian group $A$ to the sheaf defined by $\F(U)=A$ if $x\in U$ (so $x^{-1}U=*$) and $0$ otherwise (so $x^{-1}U=\emptyset$), i.e.\ the skyscraper sheaf.
    And $x^{-1}\colon\sh(X)\to\Ab$ maps $\F\in\sh(X)$ to the direct colimit $\colim_{x\in U}\F(U)$, i.e.\ the stalk of $\F$ at $x$.
    So we have shown that $\bul_x\dashv x_*\bul$.

\eexam

\bexam

    For a small category $I$, we have the {\emphcolor diagonal functor} $\Delta_I\colon\A\to\A^I$ which maps $A\in\A$ to the {\emphcolor constant functor}
    $\Delta A$ mapping $i\in I$ to $A$, and $i\to j$ to $\id_A$.
    A classical result is $\colim_I\dashv\Delta_I\dashv\lim_I$ when all co/limits of shape $I$ exist.
    This is pretty immediate: $\lim_IF$ is an object $A\in\A$ along with a natural map (the limit cone) $\lambda\colon\Delta A\to F$ which is universal in a particular way (terminal in the category of cones).
    So then for $A\in\A$ and functor $F$, a map $\Delta A\to F$ defines a unique map $A\to\lim F$ commuting with the limit cone (by terminal-ness), and vice versa.

    In particular $\colim$ is right exact, $\lim$ left exact, and $\Delta$ exact.
    $\lim,\colim$ are not necessarily exact though.

\eexam

Recall that all limits can be computed via products and equalizers.
In Abelian categories this means that all limits can be computed via products and kernels.
Similarly colimits can be computed via coproducts (direct sums) and cokernels.
Thus we have the following.

\bthrm

    An Abelian category is cocomplete (all colimits exist) iff all direct sums exist, and is complete (all limits exist) iff all products exist.

\ethrm

Note that if $\D$ is a full subcategory of a cocomplete category $\C$ such that the inclusion $i\colon\D\to\C$ has a left adjoint ($\D$ is a {\emphcolor reflective subcategory}),
$\D$ is also cocomplete.
Indeed, let $s\colon\C\to\D$ be the left adjoint then the composition $\D\xtos i\C\xto\Delta\C^I$ is just the diagonal functor on $\D$, and since we can compose adjunctions,
it has a left adjoint $s\circ\colim$.
So $s\circ\colim$ is the colimit functor of $\D$, meaning $\D$ is cocomplete.
We can also show that if $\C$ is complete, so is $\D$.

In particular it is simple to show that $\Psh(X)$ is co/complete, and since $\sh(X)$ is a reflective subcategory, it too is co/complete.
Thus we have the following co/complete categories:
$$ \Mod_R,\qquad \Psh(X),\qquad \sh(X) $$

\bdefn

    An Abelian category $\A$ satisfies axiom {\emphcolor (AB4)} if it is cocomplete and direct sums of monics are monic.

\edefn

The direct sum of maps $f_i\colon A_i\to B_i$ is the induced map $\oplus_if_i\colon\oplus_iA_i\to\oplus_iB_i$ defined by composing $f_i$ with the inclusion $B_i\to\oplus_iB_i$,
and then utilizing the universal property of direct sums.
That is, $\oplus_if_i$ is the unique map where composing with inclusion $A_i\to\oplus A_i$ is equal to $f_i$ composed with inclusion $B_i\to\oplus B_i$.

Axiom (AB4) is equivalent to saying that direct sums are exact.
Since direct sums are colimits, they are automatically right exact.
The direct sums of monics being monic is equivalent to saying that direct sums preserve monomorphisms: direct sums are thus exact.

Since homology commutes with exact functors we have the following.

\bthrm

    Suppose $\A$ satisfies (AB4) and has enough projectives.
    Let $F\colon\A\to\B$ be left adjoint, then for any collection $\set{A_i}$ in $\A$, we have
    $$ L_\bul F\parens{\bigoplus_iA_i} \cong \bigoplus_iL_\bul F(A_i) $$

\ethrm

\bproof

    The direct sum of projectives is projective ($\oplus P_i\to C$ factors as $P_i\to C$ which then lifts to $P_i\to B$, giving us $\oplus P_i\to B$).
    Since direct sums are exact, if $P_i\to A_i$ are projective resolutions then so is $\oplus P_i\to A_i$.
    Then since $F$ is left adjoint and thus commutes with colimits, in particular direct sums, we have
    $$ L_\bul F\parens{\bigoplus_iA_i} = H_\bul(F(\oplus_iP_i) \cong H_\bul(\oplus_iF(P_i)) \cong \bigoplus_iH_\bul(F(P_i)) = \bigoplus_iL_\bul F(A_i) $$
    as desired.
    \qed

\eproof

We can define the dual axiom of AB4:

\bdefn

    An Abelian category $\A$ satisfies axiom {\emphcolor (AB4\star)} if the product of epimorphisms is epic.

\edefn

$\A$ satisfies (AB4\star) iff $\A^\op$ satisfies (AB4) clearly.
Thus by duality we have:

\bcoro

    Suppose $\A$ satisfies (AB4\star) and has enough injectives.
    Let $F\colon\A\to\B$ be right adjoint, then for any $\set{A_i}$ in $\A$, we have
    $$ R_\bul F\parens{\prod_iA_i} \cong \prod_i R_\bul F(A_i) $$

\ecoro

It is not hard to see that $\Mod_R$ satisfies both (AB4) and (AB4\star), thus we have

\bcoro

    In $\Mod_R$:
    $$ \tor_\bul\parens{A,\bigoplus_iB_i}\cong \bigoplus_i\tor_\bul(A,B_i),\qquad \ext_\bul\parens{A,\prod_iB_i} \cong \prod_i\ext_\bul(A,B_i) $$

\ecoro

\bdefn

    A nonempty category $I$ is {\emphcolor filtered} if
    \benum
        \item for every $i,j\in I$, there is some $k\in I$ with arrows $i\to k\ot j$.
        \item for every two parallel arrows $u,v\colon i\tto j$ there is an arrow $w\colon j\to k$ such that $wu=wv$.
    \eenum

    A {\emphcolor filtered colimit} is the colimit of a functor $I\to\A$ where $I$ is filtered.
    A filtered colimit is denoted by $\coflim$.

\edefn

We can dualize this to {\emphcolor cofiltered} categories and {\emphcolor cofiltered limits}.
Note that posets always satisfy (2) (since in such a case $u=v$), so a filtered poset is just one where every two elements can be bound.
Filtered posets are called {\emphcolor directed}.
Filtered colimits from directed posets are often called directed limits, but for clarity I will call them directed {\it colimits}.

\blemm

    Let $I$ be filtered, and $A\colon I\to\Mod_R$ a functor.
    \benum
        \item Every element $a\in\coflim(A_i)$ is the image of some element $a_i\in A_i$ under the canonical map $A_i\to\coflim A_i$.
        \item For every $i\in I$, the kernel of the canonical map $A_i\to\coflim A_i$ is the union of kernels of the maps $\phi\colon A_i\to A_j$ (where $\phi\colon i\to j$).
    \eenum

\elemm

I explain quickly what we mean by ``canonical map''.
Recall that a colimit is given by a cone $\mu\colon A\to\Delta(\colim A)$, and thus we have maps $\mu_i\colon A_i\to\colim A$; these are the canonical maps.

\bproof

    Recall that in general the colimit of $A$ is given by the coequalizer diagram
    $$ \bigoplus_{\phi\colon i\to j}A_i \xvarrightarrows{\ f\ }[\ g\ ]\bigoplus_{i\in I}A_i \to \coflim A $$
    where $f=\bigoplus_\phi f_\phi$ and $g=\bigoplus_\phi g_\phi$ for $f_\phi,g_\phi\colon A_i\to\bigoplus_iA_i$ by setting $f_\phi$ to be the inclusion of $A_i$,
    and $g_\phi$ to be the composition of the inclusion of $A_j$ with $\phi$.
    I.e.\ %
    $$ \coflim A = \coker\parens{\bigoplus_{\phi\colon i\to j}\lambda_i-\lambda_j\phi} $$
    where $\lambda_i\colon A_i\to\bigoplus_iA_i$ is inclusion.
    The canonical maps $A_i\to\coflim A_i$ are the compositions of $\lambda_i$ with the coequalizer arrow.

    Thus $\coflim A$ is the quotient of $\bigoplus_iA_i$ by the image of $\bigoplus_{\phi\colon i\to j}\lambda_i-\lambda_j\phi$.
    Elements of $\bigoplus_iA_i$ are of the form $\sum_{j\in J}\lambda_j(a_j)$ for a finite set $J\subseteq I$.
    Since $I$ is filtered, $J$ has a bound, let it be $J\leq i$.
    Let $\phi_i\colon j\to i$ for $j\in J$ then we have that $\lambda_j=\lambda_i\phi$ in the filtered colimit (we quotient by these direct sums), so we have
    $$ \sum_{j\in J}\lambda_j(a_j) = \sum_{j\in J}\lambda_i(\phi a_j) = \lambda_i\parens{\sum_j\phi a_j} $$
    proving our first point.

    If $a_i\in A_i$ vanishes in $\coflim A_i$ then this means that $\lambda_i(a_i)=\sum_{j,k}\lambda_j(a_j)-\lambda_k(\phi_{jk}a_j)$ for finitely many $j,k$,
    $\phi_{jk}\colon j\to k$, and $a_j\in A_j$.
    Let $t$ be an upper bound for all these $i,j,k$, then we have
    $$ \lambda_t(\phi_{it}a_i) = \sum_{j,k}\lambda_t(\phi_{jt}a_j) - \lambda_t(\phi_{kt}\phi_{jk}a_j) $$
    for $\phi_{it}\colon i\to t$.
    We can assume that $\phi_{kt}\phi_{jk}=\phi_{jt}$, potentially by increasing $t$ (due to the second condition of filtered categories).
    Thus we have
    $$ \lambda_t(\phi_{it}a_i) = \sum_{j,k}\lambda_t(\phi_{jt}a_j)-\lambda_t(\phi_{jt}a_j) = 0 $$
    since $\lambda_t$ is injective, this means that $a_i\in\ker\phi_{it}$ as desired.
    The converse (if an element is in the union of kernels, it vanishes) is because $\lambda_i(a_i)=\lambda_j(\phi_{ij}a_i)=0$.
    \qed

\eproof

This lemma shows that we can define a surjection $\bigsqcup A_i\to\coflim A_i$, by mapping $a_i\in A_i$ to $\lambda_i(a_i)$.
Then $a_i,a_j$ define the same element in the filtered colimit iff $a_i-a_j$ is in the kernel of one of the maps.
Bounding $i,j$ by $k$, this means that $a_i-a_j$ is in the kernel of some $\phi\colon A_k\to A_{k'}$, i.e.\ $\phi(a_i)=\phi(a_j)$.
That is,
$$ \coflim A_i \cong \left.\bigsqcup A_i\middle/\sim\right.,\qquad a_i\sim a_j\hbox{ iff there exists some $i,j\to k$ such that $\phi_{ik}(a_i)=\phi_{jk}(a_j)$} $$

\bthrm

    Filtered colimits (in particular directed colimits) of $R$-modules are exact, as functors $\Mod_R^I\to\Mod_R$.

\ethrm

\bproof

    Let $\A=\Mod_R$, and $I$ be a filtered category.
    Colimits are in general right exact, and so $\coflim_I\colon\A^I\to\A$ is right exact.
    So we need only show that $\coflim_I$ preserves monics.
    So suppose $f\colon A\to B$ is monic in $\A^I$ (when $\A=\Mod_R$ this just means that every component of $f$ is monic).
    We want to show that the induced map $f_*\colon\coflim A_i\to\coflim B_i$ is monic as well.
    So suppose $a\in\coflim A_i$ vanishes in $\coflim B_i$, by the above lemma $a$ is the image of some $a_i\in A_i$.
    Thus $f_i(a_i)\in B_i$ vanishes in $\coflim B_i$, meaning there is some $\phi\colon i\to j$ such that $\phi(f_i(a_i))=f_j(\phi(a_i))=0$ in $B_j$.
    Since $f_j$ is monic, $\phi(a_i)=0$, so $a_i$ vanishes in $\coflim A_i$, i.e.\ $a=0$.
    \qed

\eproof

This leads us to the following definition:

\bdefn

    An Abelian category $\A$ satisfies the axiom {\emphcolor (AB5)} if it is cocomplete and filtered colimits are exact.
    Dually it satisfies {\emphcolor (AB5\star)} if it is complete and cofiltered limits are exact.

\edefn

$\Mod_R$ is not (AB5\star): $\Mod_R^\op$ is not (AB5).

\bexam

    Recall that for continuous $f\colon X\to Y$, the inverse image functor $f^{-1}\colon\sh(Y)\to\sh(X)$ is given by the sheafification of
    $f^{-1}\G(U)=\colim_{f(U)\subseteq V}\G(V)$.
    Sheafification is exact, so we need only show that $f^{-1}\colon\Psh(Y)\to\Psh(X)$ is exact to show that the inverse image functor on sheaves is.

    A sequence of presheaves (not sheaves) is exact iff it is exact at every open set.
    Note that the set of open sets $V$ containing $f(U)$ is directed, and thus $f^{-1}\G(U)$ is a filtered colimit, and thus exact.
    That is to say if $\G\to\G'\to\G''$ is exact, then these define an exact sequence of functors on the directed set of open subsets containing $f(U)$.
    Then taking their filtered colimit over this directed set gives us that $f^{-1}\G(U)\to f^{-1}\G'(U)\to f^{-1}\G''(U)$ is exact.
    This is true for all $U$, and thus $f^{-1}$ is exact, as desired.

\eexam

Since filtered colimits are exact, they commute with homology.
And left adjoints commute with colimits in general, so we have the following.

\bthrm

    Suppose $\A,\B$ satisfy (AB5), $\A$ has enough projectives.
    If $F\colon\A\to\B$ is left adjoint such that filtered limits of projectives in $\A$ are $F$-acyclic, then for every $A\colon I\to\A$ with $I$ filtered,
    $$ L_\bul F(\coflim A_i) \cong \coflim L_\bul F(A_i) $$

\ethrm

Note that the filtered limit of projective modules are $\bul\otimes A$-acylic for every $R$-module $A$.
Indeed, $\bul\otimes(\coflim P_i)$ is exact, as it is isomorphic to $\coflim(\bul\otimes P_i)$ which is the filtered colimit of exact functors, and is thus exact by (AB5).
This gives us the result (after balancing $\tor$, i.e.\ showing that it can be viewed as a bifunctor).
Thus we have the following corollary:

\bcoro

    For every filtered $B\colon I\to\Mod_R$ and $A\in{}_R\Mod$,
    $$ \tor_\bul(A,\coflim B_i) \cong \coflim \tor_\bul(A,B_i) $$

\ecoro

\subsection{Balancing Tor and Ext}

Recall that since tensor products and hom-sets are bifunctors, which are (left/right) exact in both their coordinates, they both provide two derived functors.
I.e.\ $L_\bul(\bul\otimes_RB),L_\bul(A\otimes_R\bul)$ for tensor products.
We claimed that $L_\bul(\bul\otimes_RB)(A)\cong L_\bul(A\otimes_R\bul)(B)$ (similar for $\ext$); we show this now.

\bdefn

    Let $P,Q$ be chain complexes of right and left $R$-modules, respectively.
    Form the bicomplex $P\otimes_RQ$ whose components are $P_p\otimes_RQ_q$, with horizontal differential $d_P\otimes1$, and vertical differential $(-1)^p\otimes d_Q$.
    This is called the {\emphcolor tensor product bicomplex}, and $\tot^\oplus(P\otimes_RQ)$ is the {\emphcolor tensor product chain complex} of $P$ and $Q$.

\edefn

This is indeed a bicomplex: $(d\otimes1)\circ(1\otimes d)$ maps $p\otimes q$ to $dp\otimes dq$, and thus these differentials form a positive bicomplex, we use the sign trick
to give us a signed bicomplex.

If $P$ is a chain complex of right $R$-modules and $A$ is a left $R$-module, then we can view $A$ as a chain complex concentrated at $0$, thus defining $P\otimes A$.
Note then that $\tot^\oplus(P\otimes_RA)$ has coefficient at index $n$, $P_n\otimes_RA$, since all other coefficients in $A$ are $0$.
The differential is $d_P\otimes1+(-1)^p\otimes d_A=d_P\otimes1$.

If we have chain maps $f\colon P\to A$ and $g\colon Q\to B$, these define chain (bicomplex) maps $f\otimes g\colon P\otimes Q\to A\otimes B$ whose components are $f_p\otimes g_q$.

\bthrm

    $L_n(A\otimes_R\bul)(B)\cong L_n(\bul\otimes_RB)(A)=\tor^R_n(A,B)$ for all $A,B$ and $n$.

\ethrm

\bproof

    Choose projective resolutions $P\xto\epsilon A$ in $\Mod_R$ and $Q\xto\eta B$ in ${}_R\Mod$.
    Consider $A,B$ as complexes concentrated at degree $0$, then we have three bicomplexes: $P\otimes Q$, $P\otimes B$, and $A\otimes Q$.
    The augmentations $\epsilon$ and $\eta$ define maps from $P\otimes Q$ to $A\otimes Q$ and $P\otimes B$ respectively.
    The {\emphcolor acyclic assembly lemma} (below) will help us show that the maps
    $$ A\otimes Q = \tot(A\otimes Q) \xot{\epsilon\otimes Q} \tot(P\otimes Q) \xto{P\otimes\eta} \tot(P\otimes B) = P\otimes B $$
    are quasi-isomorphisms, inducing the desired isomorphisms on homology:
    $$ L_\bul(A\otimes_R)(B) \xot{\ \cong\ } H_\bul(\tot(P\otimes Q)) \xtos\cong L_\bul(\otimes_RB)(A) $$
    (Recall that $L_\bul(A\otimes_R)(B)$ is computed by taking a projective resolution $Q\to B$ and then computing the homology of $A\otimes Q$.)

    Consider the bicomplex $C$ obtained by adding $A\otimes Q$ in the column $p=-1$ of $P\otimes Q$, so that $C_{-1,q}=A\otimes Q_q$, and the horizontal differential
    $C_{0,q}=P_0\otimes Q_q\to C_{-1,q}=A\otimes Q_q$ is given by $\epsilon\otimes1$.
    Now consider $\tot(C)$: its $n$th coefficient is $(A\otimes Q_{n+1})\oplus\tot(P\otimes Q)_n$.
    Its differentials are the direct sum of $\epsilon\otimes1-1\otimes d_Q$ with the differential of $\tot(P\otimes Q)$, $d^\tot$.
    So $\tot(C)[-1]$ has its $n$th coefficient $(A\otimes Q_n)\oplus\tot(P\otimes Q)_{n-1}$ with differential $-d^\tot\oplus(1\otimes d_Q-\epsilon\otimes1)$.

    Consider the mapping cone of $\epsilon\otimes1\colon\tot(P\otimes Q)\to A\otimes Q$.
    Its $n$th component is $\tot(P\otimes Q)_{n-1}\oplus(A\otimes Q_n)$, and its differential is $-d^\tot\oplus(1\otimes d_Q-\epsilon\otimes1)$.
    Thus $\tot(C)$ is isomorphic to the mapping cone of $\epsilon\otimes1$.
    So to show that $\epsilon\otimes1$ is a quasi-isomorphism (equivalently its mapping cone is acyclic), it is sufficient to show $C$ is acyclic.

    Since $\otimes Q_q$ is an exact functor, the rows of $C$ are exact, and $C$ is an upper half plane bicomplex.
    Thus the acyclic assembly lemma tells us $\tot C$ is acyclic, as desired.

    Similarly we can add $P\otimes B$ to the row $q=-1$ of $P\otimes Q$ and the resulting bicomplex's total complex is acyclic (by the assembly lemma) and isomorphic
    to the mapping cone of $P\otimes\eta$.
    \qed

\eproof

\blemm[title=Acyclic assembly lemma]

    Let $C$ be a bicomplex of $R$-modules.
    Then
    \benum
        \item $\tot^\Pi(C)$ is acylic if either of the following hold:
        \benum
            \item $C$ is an upper half-plane complex with exact columns, or
            \item $C$ is a right half-plane complex with exact rows.
        \eenum
        \item $\tot^\oplus(C)$ is acylic if either of the following hold:
        \benum
            \item $C$ is an upper half-plane complex with exact rows, or
            \item $C$ is a right half-plane complex with exact columns.
        \eenum
    \eenum

\elemm

\bproof

    Consider the transpose of $C$, $C^\top$.
    Then $\tot(C^\top)=\tot(C)$, and $C^\top$ is right half-plane iff $C$ is upper half-plane so $(1)$ and $(2)$ are equivalent, and $(3)$ and $(4)$ are equivalent.

    We claim that $(1)$ implies $(4)$: suppose $C$ is right half-plane with exact columns.
    Define $\tau_nC$ by truncating the columns of $C$ at level $n$:
    $$ (\tau_nC)_{pq} = \cases{C_{pq} & if $q>n$\cr \ker(d^v\colon C_{pn}\to C_{p,n-1}) & if $q=n$\cr 0 & else} $$
    Each $\tau_n$ is (up to translation) a first-quadrant bicomplex, and so its total complexes coincide: $\tot^\oplus(\tau_nC)=\tot^\Pi(\tau_nC)$.
    Furthermore $\tau_nC$ has exact columns (its columns are the truncations of the chain complex $C_{p,\bul}$ which has the same homology of $C_{p,\bul}$; zero).
    Assuming $(1)$, this means that $\tot(\tau_nC)$ is acylic.

    Now $\tot(\tau_nC)$ are subcomplexes of $\tot^\oplus(C)$, and their union (union of each component) is all of $\tot^\oplus(C)$.
    Thus every cycle of $\tot^\oplus(C)$ is a cycle in some $\tot(\tau_nC)$, and thus also a boundary since $\tot(\tau_nC)$ is acylic.
    This means that $\tot^\oplus(C)$ is acyclic, as desired.

    To show $(1)$, by translating $C$ left and right, it is sufficient to show that $H_0(\tot(C))=0$.
    Let $c=(\dots,c_{-p,p},\dots,c_{-1,1},c_{0,0})_p\in\prod_pC_{-p,p}=\tot(C)_0$ be a $0$-cycle.
    This means that $d^v(c_{-pp})+d^h(c_{-p+1,p-1})=0$.
    We find by induction elements $b_{-p,p+1}$ such that $d^v(b_{-p,p+1})+d^h(b_{-p+1,p})=c_{-p,p}$.
    The differential of $b$ is then $c$, meaning that $c$ is a boundary; and then $H_0(\tot(C))=0$ as desired.

    For $p=-1$ choose $b_{1,0}=0$.
    And for $p=0$, since $C_{0,-1}=0$ we must have that $d^v(c_{00})=0$, and since the zeroth column is exact there is a $b_{0,1}\in C_{0,1}$ such that $d^v(b_{0,1})=c_{00}$.
    This gives us $d^v(b_{0,1})+d^h(b_{1,0})=c_{00}$ as desired.

    Inductively, we have
    $$ d^v(c_{-p,p}-d^h(b_{-p+1,p})) = d^v(c_{-pp}) + d^hd^v(b_{-p+1,p}) = d^v(c_{-pp}) + d^h(c_{-p+1,p-1}) - d^hd^h(b_{-p+2,p-1}) = d^v(c_{-pp}) + d^h(c_{-p+1,p-1}) = 0 $$
    Since the $-p$th column is exact, $c_{-p,p}-d^h(b_{-p+1,p}$ must be in the image of $d^v$: there is a $b_{-p,p+1}\in C_{-p,p+1}$ such that
    $$ d^v(b_{-p,p+1}) = c_{-p,p} - d^h(b_{-p+1,p}) $$
    as desired.
    \qed

\eproof

\bdefn

    Let $P$ be a chain complex, and $I$ a cochain complex.
    We define their {\emphcolor hom bicomplex} to be the cochain bicomplex $\hom(P,I)$ with components $\hom(P,I)_{p,q}=\hom(P_p,I^q)$ and differentials
    $$ d^h\colon\hom(P_p,I^q)\to\hom(P_{p+1},I^q),\ f\mapsto f\circ d_P,\qquad d^v\colon\hom(P_p,I^q)\to\hom(P_p,I^{q+1}),\ f\mapsto (-1)^{p+q+1}d_Q\circ f $$
    Their {\emphcolor (total) hom cochain complex} is defined to be $\tot^\Pi(\hom(P,I))$.

\edefn

Suppose we have two chain complexes, $C$ and $D$.
Reindexing $D$ as a cochain complex, then an $n$-cycle $f$ of the hom cochain complex $\hom(C,D)$ is a collection of maps $f_p\colon C_p\to D^{n-p}=D_{p-n}$ whose differential is zero.
The differential is
$$ d(f_p) = (f_pd_C\colon C_{p+1}\to D_{p-n}) + ((-1)^{n+1}d_Df_p\colon C_p\to D_{p-n-1}) $$
and so is zero when $f_pd=(-1)^ndf_{p+1}$.
Equivalently, $n$-cycles are chain maps $C\to D[-n]$.
$n$-boundaries are null homotopic chain maps.
And so the cohomology grousp of the hom cochain complex, $H^n\hom(C,D)$, are the chain homotopy classes of chain maps $C\to D[-n]$, i.e.\ the hom-set of maps $C\to D[-n]$ in $\K$.

\bprop

    There is an internal tensor-hom adjunction, that is a natural isomorphism
    $$ \tot^\Pi\hom_\Ab(\tot^\oplus(P\otimes Q),I) \cong \tot^\Pi\hom_R(P,\tot^\Pi\hom_\Ab(Q,I)) $$
    for $P,Q$ right and left $R$-module chain complexes, and $I$ a cochain complex of Abelian groups.

\eprop

\bthrm

    For all $R$-modules and $n$, we have isomorphisms
    $$ \ext^n_R(A,B) = R^n\hom_R(A,\bul)(B) \cong R^n\hom_R(\bul,B)(A) $$

\ethrm

\bproof

    Choose a projective resolution $P\to A$ and an injective resolution $B\to I$.
    Then we have a first-quadrant double cochain complex $\hom(P,I)$.
    The augmentations induce maps $\hom(A,I)\to\hom(P,I)\ot\hom(P,B)$.
    As for tensor products, the mapping cones of $\hom(A,I)\to\tot\hom(P,I)$ and $\hom(P,B)\to\tot\hom(P,I)$ are translates of the total complex obtained by
    adding $\hom(A,I)$ and $\hom(P,B)$ to $\hom(P,I)$.
    By the acylic assembly lemma (or rather its dual for cochain bicomplexes), these are exact and thus quasi-isomorphisms.
    Taking cohomology gives us
    $$ R^\bul\hom(A,\bul)(B) = H^\bul(\hom(A,I)) \cong H^\bul\tot\hom(P,I) \cong H^\bul\hom(P,B) = R^\bul\hom(\bul,B)(A) $$
    as desired.
    \qed

\eproof

We can naturally extend these results to a broader class of derived functors.

\bdefn

    Let $T$ be a multivariate functor $\A_1\times\cdots\times\A_n\to\c B$ between Abelian categories such that $T$ is left exact (equivalently left exact in each of
    its components).
    $T$ is then {\emphcolor right balanced} if whenever any of its components are held constant at an injective object, $T$ becomes exact (in each of its other components).

    Dually a right exact functor is {\emphcolor left balanced} if it becomes exact when its components are held constant at a projective object.

\edefn

For example $\hom(\bul,\bul)$ is right balanced (since injectives in $\A^\op$ are projectives in $\A$), and $\bul\otimes\bul$ are left balanced.

If $T$ is right balanced, then all of its right derived functors are isomorphic in the sense that $R^\bul T(A_1,\dots,\hat A_i,\dots,A_n)(A_i)$ are all isomorphic ($\hat A_i$ means
the omission of $A_i$, replaced by an indeterminant $\bul$).
Thus for $T$ right balanced, we can construct its {\emphcolor right derived functors} $R^\bul T\colon\A_1\times\cdots\times A_n\to\c B$.
Dually for $T$ left balanced.

The proof proceeds as above.
It is sufficient to consider only bifunctors (since the rest of the variables are held constant) $T\colon\A_1\times\A_2\to\c B$.
For cochain complexes $P^\bul,Q^\bul$ we can construct a (positive) cochain bicomplex $T(P^p,Q^q)$ with horizontal and vertical differentials as expected:
$$ d^h=T(d_P,Q^\bul)\colon T(P^p,Q^q)\to T(P^{p+1},Q^q),\qquad d^v=T(P^\bul,d_Q)\colon T(P^p,Q^q)\to T(P^p,Q^{q+1}) $$

Let $A_i\in\A_i$ and injective resolutions $A_i\to I^i_\bul$, then we can consider the bicomplex $T(I^1_p,I^2_q)$.
Appending $T(A_1,I^2_\bul)$ and $T(I^1_\bul,A_2)$ give us bicomplexes whose total complexes are isomorphic to the mapping cones $T(A_1,I^2_\bul),T(I^1_\bul,A_2)\to\tot(T(I^1_\bul,T^2_\bul))$.
By the acylic assembly lemma, these bicomplexes are acylic and thus these maps are quasi-isomorphisms.
In particular the cohomology of $T(A_1,I^2_\bul)$ and $T(I^1_\bul,A_2)$, which are $R^\bul T(A_1,\bul)(A_2)$ and $R^\bul T(\bul,A_2)(A_1)$, are isomorphic as desired.

